\chapter{Biocompatibility and medical devices}
\label{vol06:ch12}

\epigraph{The biocompatibility of a long-term implantable device
refers to the ability of the device to perform its intended function
\ldots\ with the desired degree of incorporation in the host,
without eliciting any undesirable local or systemic effects in that
host.}{David F.\ Williams, \emph{Essential Biomaterials Science},
2014 \cite{text:v6c12-williams-biocompatibility-2014}.}

\chapmeta{Half-life: foundational. Archetypes: failure, ethics, interface. Exercise target: 25. Project track: read a real device 510(k) submission and identify hidden engineering choices (hybrid: document analysis plus engineering reconstruction). Hazard class: low (no patient contact, no laboratory work; the project is desk research). Reader path default: standard, with sections explicitly tagged. Named cases: Therac-25 (acc:therac-25-1985-1987); Bjork-Shiley convexo-concave heart valve (acc:bjork-shiley-cc-1979-1986); DePuy ASR metal-on-metal hip implant (acc:depuy-asr-hip-2010).}

\noindent
A medical device is a designed object that must work inside a body
it did not evolve to inhabit and that did not evolve to receive it.
The reader who finishes this chapter can distinguish biocompatibility
as a system property from biocompatibility as a material datasheet;
can name the principal failure modes of long-term implants by
mechanical, material, biological, and operational class; can sketch
the foreign-body response timeline and the design constraints it
imposes; can describe drug-device combinations as boundary-value
problems in diffusion and dosing; can read a sterilisation cycle
against the materials it touches; can navigate the FDA classification
and 510(k) versus PMA pathways and their European Union counterparts;
and can write a short forensic reconstruction of any of the three
canonical failure cases that organise this chapter.

\section{What biocompatibility means in practice}\pathtag{core}
% Figure: host response timeline from implantation to chronic state.
% Vol VI Ch 12 sec 12.2.
\begin{figure}[htbp]
\centering
\begin{tikzpicture}[
  phase/.style={draw=engblack, minimum height=1.0cm, align=center,
    font=\small},
  label/.style={font=\scriptsize\itshape, anchor=south, align=center},
]

% Time axis
\draw[->, thick] (-0.3, 0) -- (14.5, 0);
\node[font=\scriptsize] at (14.7, -0.3) {time (log)};

% Tick marks and labels
\foreach \x/\t in {0.5/0, 2.5/seconds, 5.0/hours, 7.5/days, 10.0/weeks, 12.5/months--years} {
  \draw (\x, -0.05) -- (\x, 0.05);
  \node[font=\tiny, anchor=north] at (\x, -0.05) {\t};
}

% Phase boxes
\node[phase, fill=engblue!10, minimum width=2.5cm] (adsorb)
  at (2.5, 1.4) {\textbf{Protein adsorption}\\
  (Vroman effect)};
\node[phase, fill=englime!10, minimum width=2.5cm] (acute)
  at (5.5, 1.4) {\textbf{Acute inflammation}\\
  neutrophils};
\node[phase, fill=engred!10, minimum width=2.5cm] (chronic)
  at (8.5, 1.4) {\textbf{Chronic inflammation}\\
  macrophages};
\node[phase, fill=engred!15, minimum width=2.5cm] (granuloma)
  at (11.5, 1.4) {\textbf{Foreign-body reaction}\\
  giant cells, fibrosis};

% Encapsulation (longer arrow extending right)
\node[phase, fill=engblack!10, minimum width=3cm] (capsule)
  at (12.0, 3.0) {\textbf{Fibrous capsule}\\
  (mature; weeks to years)};

% Connecting arrows
\draw[->, thick] (adsorb) -- (acute);
\draw[->, thick] (acute) -- (chronic);
\draw[->, thick] (chronic) -- (granuloma);
\draw[->, thick] (granuloma) -- (capsule);

% Side note: integration paths
\node[label, fill=engblue!5, draw=engblack, minimum width=4cm,
  minimum height=1.5cm, align=center] at (3.0, 3.5)
  {\textbf{Alternative outcomes}\\
  Osseointegration (bone)\\
  Endothelialisation (vasculature)\\
  Degradation (resorbable)};

\end{tikzpicture}
\caption[Host-response timeline]{The canonical foreign-body
response timeline at a biomaterial surface, on a logarithmic time
axis spanning seconds (initial protein adsorption) through months
to years (mature fibrous capsule). Protein adsorption is rapid
and the adsorbed layer is what later cells recognise, not the
nominal material surface; the Vroman effect describes the
sequential displacement of early-adsorbed abundant proteins
(albumin) by lower-abundance higher-affinity proteins (fibrinogen,
factor XII). Acute neutrophilic inflammation peaks within hours;
chronic macrophage-dominated inflammation persists if the implant
is not resorbed; macrophage fusion to foreign-body giant cells
and the surrounding fibrous capsule are the chronic equilibrium
for a non-degradable solid implant. The side box names the
favourable integrative paths available to specific tissue and
device types
\cite{paper:v6c12-anderson-foreign-body-2008,
paper:v6c12-vroman-effect-1969,
text:v6c12-ratner-biomaterials-2020}.}
\label{fig:vol06:ch12:host-response-timeline}
\end{figure}


A common reading treats biocompatibility as a property of a
material. A stainless-steel data sheet declares 316L
``biocompatible''; a polymer brochure marks a polyurethane formulation
as ``implant grade.'' The reading is wrong in a way that matters.
A material is biocompatible only with respect to a specific contact
type, a specific contact duration, a specific anatomical location,
a specific manufacturing history, and a specific patient population.
The same 316L bar stock is biocompatible as a Class II surgical
retractor (limited contact with intact skin) and not biocompatible
as an unprotected long-term articulating bearing in a young active
patient. The material did not change. The context did.

\engterm{Biocompatibility} is therefore the property of a
\emph{system} comprising the material, the device geometry, the
manufacturing and sterilisation process, the surgical implantation
or use procedure, the host, and the duration over which the system
is held in contact. Williams's 2014 definition makes the system
framing explicit: biocompatibility is the ability of the device to
perform its intended function with the desired degree of host
incorporation and without eliciting undesirable local or systemic
effects in that host
\cite{text:v6c12-williams-biocompatibility-2014}.

\subsection{Contact type and contact duration}

The international standard for biological evaluation of medical
devices, \engterm{ISO 10993-1}, categorises medical devices along
two axes: contact type and contact duration
\cite{std:v6c12-iso-10993-1-2018}. The categorisation tells the
manufacturer which biological-evaluation tests apply.

\engterm{Contact type} has three categories:

\begin{itemize}[leftmargin=*]
\item \emph{Surface devices} contact intact skin, mucous membranes,
  or breached or compromised surfaces. Examples: bandages, contact
  lenses, urinary catheters, dental impressions.
\item \emph{External-communicating devices} contact circulating
  blood, indirect blood paths, or other tissues. Examples: extra-
  corporeal blood-purification circuits, dialysis lines, surgical
  needles, dental restoratives.
\item \emph{Implant devices} are introduced inside the body with the
  intention of remaining for a specified duration. Examples: heart
  valves, vascular grafts, orthopaedic prostheses, intraocular
  lenses, implanted neurostimulators.
\end{itemize}

\engterm{Contact duration} has three categories:

\begin{itemize}[leftmargin=*]
\item \emph{Limited} ($\le 24$ hours of cumulative contact).
\item \emph{Prolonged} (24 hours to 30 days).
\item \emph{Long-term} (greater than 30 days).
\end{itemize}

The product of the two categorisations is a matrix that selects a
subset of the ten-plus-part ISO 10993 evaluation programme. A
short-term surface device requires cytotoxicity, sensitisation, and
irritation tests. A long-term blood-contacting implant requires the
short-term tests plus systemic toxicity, subacute toxicity,
genotoxicity, implantation, hemocompatibility, chronic toxicity,
and carcinogenicity tests
\cite{std:v6c12-iso-10993-1-2018,
method:v6c12-iso-10993-test-matrix-2020}. The matrix is not a
spreadsheet of compliance checkmarks; each test demands a study
protocol, an animal or in vitro model, a clinical pathologist's
read, and an integrated risk assessment.

A reproduction of the ISO 10993-1 matrix, suitable for orientation
to the test programme, is included as
\texttt{data/iso\_10993\_test\_matrix.csv}; the official matrix is in
the standard itself.

\subsection{What the tests measure, and what they do not}

Each ISO 10993 test addresses one identifiable mode of harm.

\engterm{Cytotoxicity} (ISO 10993-5) measures the in vitro
viability and proliferation of a standardised cell line (most often
the L929 mouse-fibroblast line) in contact with the test article,
its extract, or its conditioned medium
\cite{std:v6c12-iso-10993-5-2009}. A failing extract suggests that
a leachable component of the device kills or impairs cells. The
test is cheap, fast, and reproducible. It is also blind to host
responses that depend on intact tissue: a cytotoxic extract that
fails an L929 test could still be acceptable inside a fibrous
capsule that walls it off; a non-cytotoxic extract could still
trigger immune sensitisation that the L929 test does not see.

\engterm{Sensitisation} measures the propensity of an extract to
elicit a delayed-type hypersensitivity reaction in a guinea-pig
maximisation test or the murine local lymph node assay (LLNA).
The test catches small-molecule sensitisers (nickel, certain epoxy
amines, residual EtO breakdown products). It does not catch
large-molecule allergens whose presentation requires
human-specific MHC.

\engterm{Irritation} measures local tissue reactivity (erythema,
oedema, vesiculation) at the device contact site, scored against
a standardised rubric.

\engterm{Systemic toxicity, subacute toxicity, chronic toxicity}
assess the systemic load delivered by leachables and degradation
products. The tests run over hours, days to weeks, and weeks to
months respectively; each terminates in histopathology of major
organ systems.

\engterm{Genotoxicity} (Ames, micronucleus, chromosomal aberration)
screens leachables for mutagenic activity in standard bacterial and
mammalian-cell assays. The screen catches strongly mutagenic
small molecules; it is the weakest part of the ISO 10993 programme
for distinguishing weakly mutagenic from non-mutagenic leachables.

\engterm{Implantation} (ISO 10993-6) is the central test for
long-term implants: the device or a representative coupon is
implanted into a standardised muscle, subcutaneous, or bone site
in a small mammal, and the surrounding tissue is harvested at
fixed time points (one week, four weeks, twelve weeks, twenty-six
weeks) for histological scoring against a published rubric
\cite{std:v6c12-iso-10993-6-2016}. The rubric scores acute
inflammation, chronic inflammation, fibrosis, encapsulation,
neovascularisation, and presence of foreign-body giant cells.

\engterm{Hemocompatibility} (ISO 10993-4) measures thrombogenicity,
complement activation, hemolysis, and platelet adhesion at the
device-blood interface. The test programme is essential for
blood-contacting devices and is the most context-dependent part of
ISO 10993: a blood-contact test in a static cuvette poorly predicts
a flowing-blood implant; a flowing-blood loop in heparinised whole
blood poorly predicts an unheparinised implant.

The tests share a common epistemic structure: each is a
\emph{model} of harm in an animal or cell system, validated to
the extent that historical positive and negative controls fall
where they are expected to fall. The model is not the patient.
A device that passes the full ISO 10993 programme has been
screened against a catalogue of mechanisms by which biocompatibility
most often fails. Remaining mechanisms are caught, if at all, in
clinical investigation under ISO 14155 and in post-market
surveillance under the regulator's adverse-event reporting
requirements
\cite{std:v6c12-iso-14155-2020}.

\begin{archetype}[Failure: the host is the other half of the device]
The reader who treats biocompatibility as a material property has
designed half a device. The other half is the host: its protein
adsorption sequence, its immune cells, its anatomical placement,
its activity load, its co-medications, and its lifetime. A device
that ``passes ISO 10993'' has been screened against catalogued
failure modes in catalogued hosts. The device that fails in the
field fails in a host the catalogue did not enumerate, with a
mechanism the catalogue did not screen, on a timeline the
catalogue did not run. Biocompatibility, like reliability, is a
property of the system in service.
\end{archetype}

\subsection{Acceptable host response is not zero response}

A second common reading treats biocompatibility as the absence of
host response. The reading is also wrong in a way that matters.
Every implant elicits some host response. Surgical placement breaks
tissue; protein adsorbs onto the surface within milliseconds;
neutrophils arrive within hours; macrophages arrive within days;
and the tissue surrounding a non-degradable solid implant
matures into a fibrous capsule of one to several millimetres
thickness over weeks to months
\cite{paper:v6c12-anderson-foreign-body-2008}. The fibrous capsule
is a host response, and it is also part of how a hip prosthesis
remains stable in the bone bed, how a pacemaker can lie quietly in
a sub-clavicular pocket for a decade, and how a breast implant can
remain in place for the lifetime of the patient. The design
question is how to shape the host response, not how to suppress it.

There is even an integrative response that the engineer
deliberately courts. \engterm{Osseointegration}, the direct apposition
of mineralised bone onto a titanium-implant surface without an
intervening fibrous layer, is the favourable equilibrium that
makes a dental implant or a femoral stem succeed
\cite{paper:v6c12-park-titanium-osseointegration-2009}. The
favourable equilibrium is produced by a specific combination of
surface chemistry (titanium dioxide layer of a few nanometres),
surface roughness (one to ten micrometres, gritblasted or
acid-etched), implant geometry (mechanical retention sufficient
for primary stability), and loading history (a quiet healing
window followed by progressive load).
\engterm{Endothelialisation}, the colonisation of a vascular
implant surface by an endothelial monolayer that participates in
normal blood-vessel function, is the favourable equilibrium that
makes a vascular graft or a drug-eluting stent succeed. In both
cases the design is for a particular kind of host response.

\subsection{Three working definitions}

For chapter use, we treat three working senses of biocompatibility:

\begin{itemize}[leftmargin=*]
\item \emph{Inertness}: the material releases no biologically
  active leachables above a threshold that triggers acute toxicity,
  sensitisation, or local irritation. The bulk of ISO 10993
  evaluates this sense.
\item \emph{Tolerance}: the device achieves a stable mature
  host-response equilibrium (typically the fibrous capsule for a
  solid implant) without acute or chronic inflammation that
  threatens function.
\item \emph{Integration}: the device achieves a tissue continuity
  with the host that participates in normal physiological function
  (osseointegration, endothelialisation, neural integration,
  tissue-engineered constructs).
\end{itemize}

A device is biocompatible to the extent that it achieves the
intended sense at the intended location for the intended duration.
``Biocompatible'' alone, without qualification, is marketing.

\begin{principle}[Biocompatibility is a system property]
Biocompatibility is not a property of a material, a coating, a
sterilisation cycle, or a host. It is a property of the
configuration: material plus geometry plus manufacturing history
plus sterilisation plus implantation procedure plus host plus
duration. A change to any element is, prima facie, a change to
the biocompatibility account, and requires a re-evaluation
proportional to the change.
\end{principle}

\section{Material-tissue interactions}\pathtag{core}

The host response to a solid implant follows a sequence that is
common across materials, with material- and surface-specific
variations on the same theme. The sequence is summarised in
Figure~\ref{fig:vol06:ch12:host-response-timeline} and described
in detail by Anderson and others
\cite{paper:v6c12-anderson-foreign-body-2008}.

\subsection{The protein adsorption layer}

Within milliseconds of implant placement, plasma and tissue
proteins adsorb to the implant surface. The adsorbed monolayer is
not the bulk plasma composition. Earlier-arriving abundant proteins
(albumin, IgG) initially dominate, and are subsequently displaced
by later-arriving lower-abundance higher-affinity proteins. The
displacement sequence is the \engterm{Vroman effect}, reported by
Vroman and Adams in 1969 using recording ellipsometry on solid-
liquid interfaces \cite{paper:v6c12-vroman-effect-1969}. On
hydrophobic surfaces the displacement chain proceeds approximately
albumin $\to$ IgG $\to$ fibrinogen $\to$ high-molecular-weight
kininogen $\to$ factor XII, on timescales of seconds to minutes.
The final layer is the surface that subsequent cells recognise.

Two practical consequences follow.

First, the cell that contacts the implant does not contact the
material. It contacts the protein layer that the material adsorbed.
A small change in surface chemistry (a fluoropolymer coating, a
poly(ethylene oxide) brush, a phosphorylcholine biomimetic surface,
a heparin-bonded layer) can produce a large change in the adsorbed
protein composition and therefore a large change in the cell
response that follows. The engineer's surface-modification options
are most usefully evaluated by what they adsorb, not by what they
nominally repel.

Second, an adsorbed protein can present an integrin-binding motif
that the protein in solution does not. \engterm{Fibrinogen} is the
canonical example. Adsorbed fibrinogen exposes the gamma-chain
$\gamma_{400-411}$ sequence and the platelet binding site, which
recruit platelets and macrophages via $\alpha_{IIb}\beta_3$ and
$\alpha_M\beta_2$ integrins respectively
\cite{paper:v6c12-mosesson-fibrinogen-2005}. The fibrinogen surface
recognises ``the body'' as a thrombogenic event, even though the
implant geometry is not a wound. The downstream consequence is
acute platelet adhesion, complement activation, and the start of
the foreign-body cascade.

\subsection{Acute inflammation}

Within hours of implant placement, neutrophils migrate to the
implant surface, driven by complement activation, fibrinogen
binding to $\alpha_M\beta_2$ on neutrophils, and surgical-trauma
chemotactic signals. Neutrophils attempt to phagocytose the implant,
fail (the implant is many times larger than any phagocyte), release
proteases and reactive-oxygen species into the local environment,
and apoptose. The acute neutrophilic phase lasts hours to a day
or two; it resolves either into a return to baseline (in the rare
case where the implant fully integrates with host tissue) or into
a chronic macrophage-dominated phase.

\subsection{Chronic inflammation and the foreign-body reaction}

Macrophages arrive within days of implantation and remain at the
surface for weeks to months or, on a long-term implant,
indefinitely. A macrophage that adheres to a foreign-body surface
and is unable to phagocytose it adopts a chronic activated state.
Repeated frustrated phagocytosis stimulates the macrophage to fuse
with neighbouring macrophages, producing multinucleated
\engterm{foreign-body giant cells}. The giant cells are the
canonical histological signature of the chronic foreign-body
reaction; they secrete cytokines, growth factors, and matrix
metalloproteinases that recruit fibroblasts and remodel the local
extracellular matrix.

The macrophage-dominated state is the most dangerous chronic state
of a non-degradable implant. Persistent macrophage activation
produces:

\begin{itemize}[leftmargin=*]
\item Continuous low-grade inflammation around the implant.
\item Continuous low-grade oxidative environment, which accelerates
  degradation of polymer surfaces and corrosion of metal surfaces.
\item Continuous recruitment of fibroblasts and collagen deposition,
  producing the fibrous capsule.
\item Continuous tissue remodelling, which on bone surfaces
  produces resorption (peri-implant osteolysis) and on soft tissues
  produces contracture.
\end{itemize}

The macrophage-dominated state ends in one of three ways: the
implant is removed, the implant is encapsulated by mature fibrous
tissue that walls off the macrophage population, or the host fails
(local extrusion, infection, systemic effects from leachables).

\subsection{The fibrous capsule}

For most long-term solid implants in soft tissue, the mature
equilibrium is a fibrous capsule: a layer of dense type-I collagen
surrounding the implant, of thickness one to several millimetres,
relatively avascular, and populated by quiescent fibroblasts and
a thin macrophage layer at the implant interface. The fibrous
capsule isolates the implant from the host. It is the favourable
chronic state for many devices: a breast implant, a cardiac
pacemaker, a glucose sensor, a sub-cutaneous infusion port.

The capsule is not always favourable. \engterm{Capsular
contracture} is the pathological tightening of the fibrous capsule
around a breast implant, producing distortion, pain, and sometimes
need for surgical revision. The contracture is graded I to IV
(Baker scale); grade IV requires capsulectomy. The mechanism is
not single. Persistent low-grade infection (biofilm), surface
texture, leakage of silicone, and surgical technique all
contribute.

The capsule is also a barrier. For a drug-eluting subcutaneous
implant, the diffusion barrier of a one-millimetre collagen-dense
capsule shifts the effective release kinetics. For an implanted
glucose sensor, capsule formation reduces sensor accuracy over
weeks. For an implanted neurostimulator electrode, capsule
formation increases impedance and demands larger drive currents.
The engineer who treats the capsule as ``just a healing response''
misses an integral part of the device's long-term function.

\subsection{Integrative responses: osseointegration, endothelialisation, neural integration}

A subset of devices courts a different equilibrium. Where the
designer wants tissue to grow into or onto the implant surface
without an intervening fibrous layer, the design must defeat the
default fibrous-capsule outcome.

\engterm{Osseointegration} is the direct functional and structural
connection between living bone and a load-bearing implant surface.
The canonical case is the titanium dental implant or the femoral
stem of a cementless total hip replacement. The conditions for
osseointegration are well characterised: a passivated titanium
oxide surface (TiO$_2$ of approximately 5 nm thickness, formed
spontaneously in air and stable in physiological saline), a
controlled surface roughness ($S_a$ in the 1 to 2 $\mu$m range,
produced by gritblasting and acid etching), primary mechanical
stability (micromotion below ~150 $\mu$m during healing), and a
quiet healing window of weeks to months before progressive
loading. Under these conditions, osteoblasts attach to the
titanium surface, deposit collagen and mineral matrix, and form
a tissue interface that approaches the load-bearing capacity of
the surrounding cortical bone
\cite{paper:v6c12-park-titanium-osseointegration-2009}.

The conditions are restrictive. A roughened titanium surface
contaminated with hydrocarbons by exposure to air at room
temperature for hours to days re-introduces the fibrous-encapsulation
default. Excess micromotion during healing produces a fibrous
interface that mechanically isolates the implant. Infection
during healing terminates the integration. The osseointegration
success rate, in well-controlled clinical conditions, exceeds 95
percent at five years for the canonical dental implant; under
field conditions it is lower.

\engterm{Endothelialisation} is the favourable equilibrium for
blood-contacting devices. A mature endothelial monolayer on the
luminal surface of a vascular graft, a heart valve, or a stent
restores the thromboresistant and inflammation-modulating function
that a healthy native vessel provides. The conditions for
endothelialisation are tighter than for osseointegration: the
surface must support endothelial-cell adhesion (often through a
collagen or fibronectin precoat), the surface must be sufficiently
permissive to endothelial migration from the anastomosis or from
the adjacent vessel, and the local fluid mechanics must not
chronically exceed endothelial shear-tolerance limits. The success
rate of endothelialisation on synthetic vascular grafts has
remained the principal limit on small-diameter graft performance
for fifty years.

\engterm{Neural integration} is the open frontier. A chronic
neural-recording or neural-stimulation electrode aims for an
interface in which neurons remain alive and electrically accessible
at distances of micrometres from the electrode. The default
response is a glial scar (the central-nervous-system analogue of
the fibrous capsule) that displaces neurons by tens of micrometres
and progressively increases electrode impedance over months. The
emerging design responses are smaller electrode geometries that
reduce the volumetric disturbance to neurons, softer materials
that match the modulus of brain tissue, bioactive coatings that
reduce glial scarring, and adaptive electronics that track
impedance over time.

\section{Implant failure modes}\pathtag{core}
% Figure: implant failure-mode tree.
% Vol VI Ch 12 sec 12.3.
\begin{figure}[htbp]
\centering
\begin{tikzpicture}[
  root/.style={draw=engblack, fill=engred!15, minimum height=0.9cm,
    minimum width=4cm, align=center, font=\small\bfseries},
  cat/.style={draw=engblack, fill=engblue!10, minimum height=0.8cm,
    minimum width=2.8cm, align=center, font=\small},
  sub/.style={draw=engblack, fill=engblack!5, minimum height=0.7cm,
    minimum width=2.6cm, align=center, font=\scriptsize},
  arrow/.style={->, thick, engblack},
]

% Root
\node[root] (root) at (0, 6) {Implant failure};

% Categories
\node[cat] (mech)  at (-6, 4.4) {Mechanical};
\node[cat] (mat)   at (-2, 4.4) {Material};
\node[cat] (bio)   at ( 2, 4.4) {Biological};
\node[cat] (op)    at ( 6, 4.4) {Operational};

% Mechanical sub-modes
\node[sub] (fat)  at (-6, 3.0) {Fatigue (cyclic)};
\node[sub] (wear) at (-6, 2.1) {Wear, fretting};
\node[sub] (lock) at (-6, 1.2) {Locking, dislocation};
\node[sub] (over) at (-6, 0.3) {Overload, impact};

% Material sub-modes
\node[sub] (corr) at (-2, 3.0) {Corrosion};
\node[sub] (creep)at (-2, 2.1) {Creep, deformation};
\node[sub] (degr) at (-2, 1.2) {Polymer degradation};
\node[sub] (em)   at (-2, 0.3) {Embrittlement};

% Biological sub-modes
\node[sub] (infl) at (2, 3.0) {Chronic inflammation};
\node[sub] (inf)  at (2, 2.1) {Infection, biofilm};
\node[sub] (ion)  at (2, 1.2) {Metal-ion toxicity};
\node[sub] (caps) at (2, 0.3) {Capsule contracture};

% Operational sub-modes
\node[sub] (sw)   at (6, 3.0) {Software fault};
\node[sub] (ui)   at (6, 2.1) {Operator interface};
\node[sub] (pos)  at (6, 1.2) {Misposition};
\node[sub] (ser)  at (6, 0.3) {Service / battery};

% Connecting lines
\draw[arrow] (root) -- (mech);
\draw[arrow] (root) -- (mat);
\draw[arrow] (root) -- (bio);
\draw[arrow] (root) -- (op);

\draw[engblack] (mech) -- (fat);
\draw[engblack] (fat) -- (wear);
\draw[engblack] (wear) -- (lock);
\draw[engblack] (lock) -- (over);

\draw[engblack] (mat) -- (corr);
\draw[engblack] (corr) -- (creep);
\draw[engblack] (creep) -- (degr);
\draw[engblack] (degr) -- (em);

\draw[engblack] (bio) -- (infl);
\draw[engblack] (infl) -- (inf);
\draw[engblack] (inf) -- (ion);
\draw[engblack] (ion) -- (caps);

\draw[engblack] (op) -- (sw);
\draw[engblack] (sw) -- (ui);
\draw[engblack] (ui) -- (pos);
\draw[engblack] (pos) -- (ser);

% Side note: coupling
\node[draw=engblack, fill=englime!10, minimum width=8cm,
  align=center, font=\scriptsize\itshape] at (0, -1.4)
  {Failure modes couple: fatigue produces wear debris that
  triggers chronic inflammation; corrosion accelerates fatigue;
  software fault and operator interface together produce
  misadministration. Single-cause boxes are pedagogical;
  released devices fail through chains.};

\end{tikzpicture}
\caption[Implant failure-mode tree]{A first-pass taxonomy of
implant failure modes by category. The four columns (mechanical,
material, biological, operational) are pedagogical; in service
they couple. Bjork-Shiley failed through fatigue (mechanical)
amplified by weld defects (material) and not caught by the
quality system (operational). DePuy ASR failed through edge
loading (mechanical) producing wear debris (mechanical) that
elevated metal-ion concentrations (material / biological) and
triggered chronic inflammation (biological); the regulatory
clearance (operational) did not anticipate the chain. Therac-25
failed through software (operational) when the migration from
hardware to software interlocks (operational) removed the
backup; the radiation overdose (biological) was the patient-
facing consequence. The single-failure-mode boxes are a useful
first taxonomy, not a complete account
\cite{text:v6c12-ratner-biomaterials-2020,
paper:v6c12-cohen-mom-bmj-2012}.}
\label{fig:vol06:ch12:failure-tree}
\end{figure}


A first-pass taxonomy of implant failure modes is the four-column
tree in Figure~\ref{fig:vol06:ch12:failure-tree}: mechanical,
material, biological, and operational. The columns are
pedagogical; in service they couple. The Bjork-Shiley failure is
mechanical (fatigue) amplified by material (weld defect) and not
caught by operational (quality system). The DePuy ASR failure is
mechanical (edge loading) producing material/biological
(metal-ion release) chains the operational (510(k)) clearance
did not anticipate. Therac-25 is operational (software) producing
biological (radiation overdose) at the patient.

\subsection{Mechanical: fatigue, wear, fretting, dislocation}

Most implants live under cyclic load. A knee implant accumulates
6,000 to 12,000 gait cycles per day, approximately
$2 \times 10^6$ to $4 \times 10^6$ cycles per year, and on a
15-year service horizon reaches $3 \times 10^7$ to $6 \times 10^7$
cycles. A heart valve sees one cardiac cycle per second on
average, approximately $3.6 \times 10^7$ cycles per year
\cite{paper:v6c12-grunkemeier-mechanical-valve-2000}. Vertebral
disc replacements, cochlear-implant lead wires, pacemaker leads,
and finger joint prostheses similarly accumulate cycle counts
that take them well into the fatigue regime over the device's
service life.

The relevant design quantity is the cyclic stress amplitude
$\Delta \sigma$ and the cycle count to crack initiation. For a
metallic alloy, the Paris-Erdogan law gives the steady-state
crack-growth rate
\[
\frac{da}{dN} = C \, (\Delta K)^m, \quad
\Delta K = Y \, \Delta \sigma \, \sqrt{\pi a},
\]
with material constants $C$ and $m$ and a geometric factor $Y$
\cite{paper:v6c12-paris-erdogan-1963}. For surgical cobalt-chromium
and titanium alloys, $m$ is approximately 3 to 4 and $C$ is
approximately $10^{-12}$ m/cycle/(MPa$\sqrt{\text{m}}$)$^m$
\cite{text:v6c12-callister-materials-2018}. The implication is
sensitive to the initial flaw size: a 50 $\mu$m flaw at a stressed
location grows at an entirely different rate than a 5 $\mu$m flaw.
Manufacturing quality of the surface and of welds is the dominant
determinant of fatigue life at population scale.

\engterm{Wear} between articulating surfaces (femoral head on
acetabular cup, knee tibial-femoral bearing, vertebral disc faces)
produces submicrometre particles that are both a mechanical and a
biological problem. The particles are below the threshold for
phagocytosis by individual macrophages, are taken up cooperatively,
and trigger chronic macrophage activation
\cite{paper:v6c12-rejman-nanoparticle-2004}. The downstream
consequence is peri-implant osteolysis: bone resorption around the
implant that progressively loosens it. The cycle-to-revision time
for a UHMWPE-on-CoCr knee or hip is set principally by the wear
rate, the particle burden, and the host's response.

\engterm{Fretting} is small-amplitude oscillatory motion between
nominally fixed surfaces (femoral head on stem taper, modular
neck-stem junction, plate-screw interface). Fretting produces wear
debris in nominally protected locations and accelerates corrosion
of passivated surfaces.

\engterm{Dislocation, loosening, locking} are macroscopic failure
modes that are easier to recognise but no less mechanical:
loss of fixation of an acetabular cup, dislocation of a femoral
head from a damaged liner, locking of a hinged knee in an unintended
angle. They are usually consequences of one of the smaller-scale
mechanisms (wear-driven loosening, capsule-driven loosening,
fatigue of a stabilising element).

\subsection{Material: corrosion, creep, degradation, embrittlement}

\engterm{Corrosion} of metallic implants in physiological saline
is the slow loss of material by electrochemical reaction. Modern
surgical alloys (316L stainless, Co-Cr-Mo, Ti-6Al-4V) are
passivated by a thin oxide layer. Under static conditions in
neutral saline, corrosion rates are negligible. Under fretting,
under mechanical surface damage, or in low-oxygen crevices
(pitting, crevice corrosion), the passive film is disrupted; local
corrosion accelerates by orders of magnitude
\cite{paper:v6c12-broensted-corrosion-stainless-2017}. The
products are not generally toxic in bulk, although the local
concentrations of cobalt and chromium can rise to biologically
active levels (compare the metal-ion release in the DePuy ASR
case, where wear and corrosion compound).

\engterm{Creep} is time-dependent deformation under static load.
UHMWPE acetabular cups creep under load over years, producing
``cold flow'' deformations that change the contact geometry of the
articulating surfaces. PMMA bone cement creeps under cyclic load,
contributing to the long-term loosening of cemented prostheses.

\engterm{Polymer degradation} encompasses oxidation of UHMWPE
(particularly after gamma sterilisation in air, which leaves free
radicals that progressively oxidise the polymer chains, embrittle
the surface, and accelerate wear), hydrolysis of bioresorbable
polymers (PLA, PGA, PLGA, PCL, designed to degrade and sometimes
with unintended kinetics), and chain scission of polyurethanes
under oxidative attack by chronic-inflammation phagocytes.

\engterm{Embrittlement} is the loss of ductility of a material
under environmental attack (hydrogen embrittlement of high-strength
steels, stress-corrosion cracking of certain stainless steels, the
embrittlement of gamma-irradiated UHMWPE).

\subsection{Biological: chronic inflammation, infection, ion toxicity, capsule contracture}

\engterm{Chronic inflammation} that fails to resolve into a stable
encapsulated state is itself a failure mode. Persistent macrophage
activation degrades polymer surfaces, drives metal corrosion,
recruits osteoclasts in bone (peri-implant osteolysis), and
shortens functional implant life.

\engterm{Infection} is the catastrophic chronic-implant
complication. A bacterial biofilm on an implant surface is
substantially shielded from systemic antibiotic concentrations
(typical biofilm minimum inhibitory concentrations are one to three
orders of magnitude above planktonic minimum inhibitory
concentrations) and from host immune cells. The standard
treatment for an infected joint prosthesis is two-stage revision:
remove the implant, place an antibiotic-loaded cement spacer, treat
the patient with systemic antibiotics for six weeks, then implant
a new prosthesis. The rate of infection for primary hip arthroplasty
is approximately 1 to 2 percent at one year, with the per-patient
lifetime risk higher for revisions. Infection is the single most
costly biological complication in the implant field and the
single largest target of surface-modification research
(antimicrobial coatings, controlled-release antibiotic systems,
copper- or silver-doped surfaces, electrochemical modulation).

\engterm{Metal-ion toxicity} from implant-derived cobalt, chromium,
nickel, and titanium has been recognised since the early metal-
on-metal hip experience and was central to the DePuy ASR failure.
Local elevation of cobalt and chromium produces ALVAL (aseptic
lymphocytic vasculitis-associated lesion), pseudotumours, and
adverse local tissue reaction. Systemic elevation (cobalt above
approximately 7 $\mu$g/L whole blood) is associated with
cardiomyopathy, neurological symptoms, and thyroid dysfunction in
case reports
\cite{paper:v6c12-jacobs-metal-ions-2003,
paper:v6c12-cohen-mom-bmj-2012}. The MHRA threshold for
investigation of MoM hip patients is currently 7 $\mu$g/L for
either cobalt or chromium
\cite{web:v6c12-mhra-mom-2012}.

\engterm{Capsule contracture}, as discussed in section 12.2, is
the chronic complication of soft-tissue implants and is the most
common reason for breast-implant revision.

\subsection{Operational: software fault, operator interface, misposition, service}

The operational column collects failure modes that depend on the
device-in-use rather than on the device-as-built. \engterm{Software
fault} (Therac-25; section 12.7) produces patient harm even with
no material or mechanical failure. \engterm{Operator-interface
mismatch} (Therac-25's MALFUNCTION 54 display; numerous infusion-
pump events) produces patient harm when the device's information
about its safety state is inaccessible to the operator at the
right moment. \engterm{Misposition} (DePuy ASR cup at steep
inclination; section 12.7) shifts the device into a regime its
design did not anticipate. \engterm{Service failures} (battery
depletion in implanted pacemakers and ICDs, lead-wire fracture in
chronic neurostimulators, dialysis-circuit set-up errors) recur
in the post-market surveillance record at rates that justify
attention.

\subsection{Coupling and failure chains}

Single failure modes do not adequately describe what kills
released devices in the field. The chain is the unit of analysis.

A cyclically loaded weld at the BSCC outlet strut accumulates
fatigue damage (mechanical); the weld contains a defect from
manufacturing (material); the quality system does not detect the
defect (operational); the strut fractures and disc embolises
(biological consequence at the patient).

A small-clearance metal-on-metal hip bearing operates with
hydrodynamic lubrication when correctly positioned (mechanical);
implanted at a steeper-than-intended inclination, edge loading
collapses the lubrication film (mechanical); metal-on-metal
contact produces submicrometre wear debris and elevated metal-ion
release (material); chronic inflammation and ALVAL ensue
(biological); the regulatory pathway did not anticipate the
chain (operational).

A radiation-therapy machine's safety interlocks are migrated from
hardware to software (operational); a race condition in the
software allows the high-energy electron beam to fire without
the metal scattering target in place (operational); the operator
interface communicates a routine malfunction code without
distinguishing dose anomaly from a recoverable beam-positioning
error (operational); the patient receives a massive overdose
(biological).

The failure chain is the design object. The chain has links in
all four columns. A design discipline that protects against
single-link failures in only one column is not protecting against
the chain.

\section{Drug-device combinations}\pathtag{standard}
% Figure: stent geometry and radial expansion.
% Vol VI Ch 12 sec 12.2 / 12.4.
\begin{figure}[htbp]
\centering
\begin{tikzpicture}[
  arrow/.style={->, thick, engblack},
  label/.style={font=\scriptsize},
]

% Crimped (delivery) configuration on the left
\node[font=\small\bfseries] at (-3.5, 3.0) {Crimped (delivery)};
\draw[engblue, very thick] (-3.5, 1.5) circle (0.5);
\draw[engblack, very thick, dashed] (-3.5, 1.5) circle (0.6);
\node[label, anchor=west] at (-3.0, 0.6) {$d_{\text{crimp}} \approx 1.0\text{--}1.5\,\text{mm}$};

% Expanded configuration on the right
\node[font=\small\bfseries] at (3.5, 3.0) {Expanded (deployed)};
\draw[engblue, very thick] (3.5, 1.5) circle (1.5);
\draw[engblack, very thick, dashed] (3.5, 1.5) circle (1.7);
\node[label, anchor=west] at (4.5, 0.6) {$d_{\text{nom}} \approx 2.5\text{--}5.0\,\text{mm}$};

% Strut pattern: simplified zigzag
\foreach \a in {0, 30, 60, 90, 120, 150, 180, 210, 240, 270, 300, 330} {
  \draw[engblue, thick] ({-3.5 + 0.5 * cos(\a)}, {1.5 + 0.5 * sin(\a)}) --
    ({-3.5 + 0.6 * cos(\a + 15)}, {1.5 + 0.6 * sin(\a + 15)});
}
\foreach \a in {0, 30, 60, 90, 120, 150, 180, 210, 240, 270, 300, 330} {
  \draw[engblue, thick] ({3.5 + 1.5 * cos(\a)}, {1.5 + 1.5 * sin(\a)}) --
    ({3.5 + 1.7 * cos(\a + 15)}, {1.5 + 1.7 * sin(\a + 15)});
}

% Arrow between
\draw[arrow, ultra thick] (-2.5, 1.5) -- (1.5, 1.5);
\node[label, anchor=south, align=center] at (-0.5, 1.6)
  {balloon inflation\\$P \sim 6\text{--}20\,\text{atm}$};

% Radial stress diagram below
\node[font=\small\bfseries, anchor=west] at (-4.5, -1.0)
  {Hoop stress balance:};
\node[font=\small, anchor=west] at (-4.5, -1.7)
  {$\sigma_{\theta} = \dfrac{P R}{t}$ (thin wall)};
\node[font=\scriptsize, anchor=west, engblack!70] at (-4.5, -2.5)
  {$P$ = balloon pressure (acting on artery wall);};
\node[font=\scriptsize, anchor=west, engblack!70] at (-4.5, -2.9)
  {$R$ = expanded inner radius;};
\node[font=\scriptsize, anchor=west, engblack!70] at (-4.5, -3.3)
  {$t$ = strut wall thickness ($60\text{--}100\,\mu\text{m}$ modern).};

% Stress-strain curve
\draw[->, thick] (3.5, -3.5) -- (6.5, -3.5);
\node[label, anchor=west] at (6.5, -3.5) {$\varepsilon$};
\draw[->, thick] (3.5, -3.5) -- (3.5, -1.0);
\node[label, anchor=east] at (3.5, -1.0) {$\sigma$};

% Stainless steel curve: elastic-plastic
\draw[engred, very thick] (3.5, -3.5) -- (4.0, -2.0)
  to[out=70, in=190] (5.0, -1.5)
  -- (6.2, -1.4);
\node[label, engred, anchor=west, align=left] at (4.0, -2.5)
  {Co-Cr or 316L\\stainless};

% Yield marker
\draw[engblack, dashed] (4.0, -3.5) -- (4.0, -2.0);
\draw[engblack, dashed] (3.5, -2.0) -- (4.0, -2.0);
\node[label, anchor=north] at (4.0, -3.5) {$\varepsilon_y$};

\end{tikzpicture}
\caption[Coronary stent geometry and expansion]{Schematic of a
coronary balloon-expandable stent in its crimped delivery
configuration (left) and expanded deployed configuration (right).
The stent is delivered on a balloon catheter at approximately
$1.0\text{--}1.5\,\text{mm}$ outer diameter through a guide
catheter, advanced to the lesion, and expanded by the balloon to a
nominal $2.5\text{--}5.0\,\text{mm}$ inner diameter at inflation
pressures of $6\text{--}20\,\text{atm}$. The strut metal (modern
cobalt-chromium or 316L stainless steel; some bioresorbable
designs use poly-L-lactide) is deformed plastically past yield;
recoil after balloon deflation is the elastic spring-back, around
3 to 7 percent for cobalt-chromium. The thin-wall hoop-stress
balance $\sigma_\theta = PR/t$ sets the strut thickness required
for a given inflation pressure and lumen radius
\cite{paper:v6c12-stone-des-2010,
text:v6c12-ratner-biomaterials-2020}.}
\label{fig:vol06:ch12:stent-geometry}
\end{figure}

% Figure: drug-eluting stent cross-section with concentration profile.
% Vol VI Ch 12 sec 12.4.
\begin{figure}[htbp]
\centering
\begin{tikzpicture}[
  layer/.style={draw=engblack, minimum height=0.5cm, align=center,
    font=\scriptsize},
  arrow/.style={->, thick, engblack},
]

% Stent strut cross-section: left half is the strut, right side is tissue
% Strut (metal)
\fill[engblack!30] (-3.5, -0.6) rectangle (0, 0.6);
\node[font=\small, white] at (-1.7, 0) {\textbf{stent strut (Co-Cr)}};

% Drug-loaded polymer coating
\fill[engblue!40] (0, -0.6) rectangle (0.6, 0.6);
\node[font=\tiny, anchor=center, rotate=90] at (0.3, 0)
  {polymer + drug};

% Tissue (arterial intima / media)
\fill[engred!10] (0.6, -0.6) rectangle (5.5, 0.6);
\node[font=\small] at (3, 0) {\textbf{arterial tissue}};

% Boundary labels
\draw[engblack, thick] (0, -0.6) -- (0, 0.6);
\draw[engblack, thick] (0.6, -0.6) -- (0.6, 0.6);

% Annotation: thickness L
\draw[<->] (0, -0.9) -- (0.6, -0.9);
\node[font=\scriptsize, anchor=north] at (0.3, -0.9)
  {$L \approx 5\text{--}15\,\mu\text{m}$};

% Concentration profile on top
\draw[->, thick] (-3.5, 1.5) -- (5.7, 1.5);
\node[font=\scriptsize, anchor=west] at (5.7, 1.5) {$x$};

\draw[->, thick] (-3.5, 1.5) -- (-3.5, 4.0);
\node[font=\scriptsize, anchor=east] at (-3.5, 4.0) {$C(x, t)$};

% Concentration profile: high in coating, decreasing into tissue
\draw[engblue, very thick] plot[domain=0:0.6, samples=10]
  (\x, {3.5 - 0.5 * \x});
\draw[engblue, very thick] plot[domain=0.6:5.5, samples=40]
  (\x, {3.2 * exp(-0.7 * (\x - 0.6)) + 1.5});

% Time progression: dashed for later time
\draw[engblue, very thick, dashed] plot[domain=0:0.6, samples=10]
  (\x, {2.4 - 0.3 * \x});
\draw[engblue, very thick, dashed] plot[domain=0.6:5.5, samples=40]
  (\x, {2.0 * exp(-0.4 * (\x - 0.6)) + 1.55});

% Labels for the curves
\node[font=\scriptsize, anchor=west, engblue] at (1.0, 3.5)
  {$t_1$ (early)};
\node[font=\scriptsize, anchor=west, engblue] at (1.5, 2.4)
  {$t_2$ (late)};

% Vertical dashed lines at boundaries
\draw[engblack!50, dashed] (0, 1.5) -- (0, 4.0);
\draw[engblack!50, dashed] (0.6, 1.5) -- (0.6, 4.0);

% Bottom: cumulative release plot
\draw[->, thick] (-3.5, -2.0) -- (5.5, -2.0);
\node[font=\scriptsize, anchor=west] at (5.5, -2.0) {$t$};
\draw[->, thick] (-3.5, -2.0) -- (-3.5, -0.7);
\node[font=\scriptsize, anchor=east] at (-3.5, -0.8)
  {$M_t / M_\infty$};

% Crank-style release curve: 1 - exp(-kt) approximation
\draw[engred, very thick] plot[domain=0:8.5, samples=40]
  (-3.5+\x, {-2.0 + 1.0 * (1 - exp(-0.4 * \x))});

\node[font=\scriptsize, anchor=west] at (-2.5, -1.3)
  {Cumulative release $\sim 1 - \exp(-k t)$ (long-time approx)};

\end{tikzpicture}
\caption[Drug-eluting stent cross-section]{Cross-section of a
drug-eluting stent strut with a polymer drug-reservoir coating of
thickness $L \approx 5\text{--}15\,\mu\text{m}$, in contact with
arterial tissue. The upper plot shows the radial drug
concentration $C(x, t)$ at an early time $t_1$ and at a late time
$t_2$: the coating depletes from a near-uniform initial loading,
and the tissue boundary layer carries the highest tissue
concentration. The lower plot shows the cumulative release
$M_t/M_\infty$ with a Crank-style one-dimensional diffusion solution;
the early-time regime is approximately Higuchi's
$\sqrt{t}$ law, the late-time regime is approximately
exponential with characteristic time $\tau \sim L^2/D$. The
design objective is the antiproliferative drug
window: too fast and the tissue concentration spikes above the
toxic threshold; too slow and the proliferative response is not
suppressed and the stent re-stenoses
\cite{paper:v6c12-crank-diffusion-1975,
paper:v6c12-higuchi-drug-release-1961,
paper:v6c12-stone-des-2010,
paper:v6c12-finn-des-pathology-2007}.}
\label{fig:vol06:ch12:des-cross-section}
\end{figure}


A \engterm{drug-device combination product} integrates a drug or
biologic with a device in a way that neither component alone could
achieve the intended therapy. The taxonomy includes:

\begin{itemize}[leftmargin=*]
\item Drug-eluting stents (rapamycin, paclitaxel, everolimus on a
  coronary stent platform).
\item Drug-eluting orthopaedic implants (antibiotic-loaded bone
  cement; vancomycin-loaded titanium implants).
\item Drug-eluting intravaginal rings, intrauterine devices,
  subcutaneous contraceptive implants.
\item Transdermal patches (nicotine, fentanyl, scopolamine).
\item Microneedle arrays for vaccines and local-effect drugs.
\item Bone graft substitutes with growth-factor delivery
  (BMP-2 on collagen sponge, for example).
\item Sustained-release ocular implants (steroid for diabetic
  macular oedema; antiviral for cytomegalovirus retinitis).
\end{itemize}

The engineering problem in each case is a boundary-value problem
in diffusion (or convection plus diffusion), with the additional
constraint that the drug concentration in the target tissue must
remain in a therapeutic window. Too low and the drug does not
suppress the proliferative response, the infection, the
inflammation, the pregnancy, or the symptom; too high and the
drug exceeds the local-toxicity threshold or the systemic-
exposure threshold.

\subsection{Diffusion from a coated reservoir}

Consider a coated coronary stent strut. The drug is dissolved or
suspended in a polymer matrix of thickness $L$ (typically 5 to 15
$\mu$m), supported on a metallic strut, in contact with the
arterial intima. The release problem is one-dimensional diffusion
into an effectively infinite sink.

For a slab of polymer with uniform initial concentration $C_0$,
diffusivity $D$, and a perfect sink at the tissue boundary, the
cumulative-release fraction follows Crank's solution
\cite{paper:v6c12-crank-diffusion-1975}:
\[
\frac{M_t}{M_\infty} = 1 - \frac{8}{\pi^2}
\sum_{n=0}^\infty \frac{1}{(2n+1)^2}
\exp\!\left[-\frac{(2n+1)^2 \pi^2 D t}{L^2}\right].
\]
At short times the leading-term approximation is Higuchi's
$\sqrt{t}$ law,
\[
\frac{M_t}{M_\infty} \approx \frac{4}{L}\sqrt{\frac{D t}{\pi}},
\]
valid until approximately 60 percent of the drug has released
\cite{paper:v6c12-higuchi-drug-release-1961}. At long times the
single-mode approximation is
\[
\frac{M_t}{M_\infty} \approx 1 - \frac{8}{\pi^2}
\exp\!\left[-\frac{\pi^2 D t}{L^2}\right],
\]
with characteristic time $\tau = L^2 / (\pi^2 D)$.

The design problem is to pick $L$, $D$, and $C_0$ to put the
tissue concentration in the therapeutic window for the intended
duration. For a sirolimus-eluting stent with drug-polymer
diffusivity of approximately $10^{-12}$ cm$^2$/s, a 10 $\mu$m
coating gives $\tau$ on the order of 30 days, sufficient to
suppress neointimal proliferation through the critical window.
For a paclitaxel-eluting stent the diffusivity is higher and the
characteristic time is shorter; the tradeoff is reflected in the
clinical-outcome differences between drug platforms
\cite{paper:v6c12-stone-des-2010}.

The model has limits that the engineer needs to be ready to
acknowledge. The polymer matrix is rarely uniform, the diffusivity
depends on drug loading and on polymer crystallinity, the tissue
boundary is not a perfect sink, and the geometry is not infinite.
Real designs use the closed-form model as a first-cut sizing tool
and confirm the prediction by in vitro elution into a buffered
medium, by ex vivo perfusion in an arterial preparation, and by
in vivo pharmacokinetics in a porcine coronary model.

\subsection{The late-stent-thrombosis problem}

Drug-eluting stents (DES) reduced restenosis rates from
approximately 25 to 35 percent at six months (the bare-metal-stent,
BMS, era) to under 5 percent (modern second-generation DES,
\cite{paper:v6c12-stone-des-2010}). The cost was an unanticipated
new failure mode: late stent thrombosis at six months to several
years post-implantation, with case-fatality rates of approximately
30 percent. The mechanism, as reconstructed by pathology series,
was the delayed endothelialisation of stent struts in the
presence of an antiproliferative drug
\cite{paper:v6c12-finn-des-pathology-2007}. The drug that suppressed
the proliferative response also suppressed the protective
endothelial response. The struts remained exposed to platelets and
to the coagulation cascade for months instead of weeks.

The design response was two-fold. First, dual antiplatelet therapy
(aspirin plus clopidogrel) was extended from three months (BMS
standard) to twelve months and longer (DES standard). Second,
second-generation DES platforms used drugs with shorter
antiproliferative duration, more biocompatible polymer coatings,
and thinner strut profiles. The combined effect was a reduction
in late-thrombosis incidence to approximately 0.2 to 0.5 percent
per year by 2010 to 2015, against the favourable restenosis-rate
benefit.

The DES history is the canonical drug-device case for the
chapter's central principle: changing one element of the system
(adding a drug to the stent) changes the host response (delayed
endothelialisation), which changes the required clinical management
(extended dual antiplatelet therapy), which feeds back into device
design (newer drugs, newer polymers, newer strut profiles). The
device is not separable from the drug, from the host, from the
co-medication, or from the time course.

\subsection{Antibiotic-loaded bone cement}

PMMA bone cement loaded with gentamicin or vancomycin is the
oldest drug-device combination at scale, in continuous use in
orthopaedic infection prophylaxis and treatment since the 1970s.
The release kinetics are imperfect (the burst release is fast and
the long tail is below the minimum inhibitory concentration of
many target organisms), the elution depends on the cement's
porosity, and the long-term retained antibiotic acts as a selection
pressure for resistant organisms. The design is sufficient for
many cases (revisions of infected total joint replacements,
spacer placements during two-stage revision) and unsuitable for
others (prophylaxis in low-infection-rate primary procedures,
where the resistance pressure outweighs the absolute benefit).

\subsection{Subcutaneous and intravaginal contraceptive implants}

Etonogestrel-releasing subcutaneous implants (the Nexplanon device,
for example) and levonorgestrel-releasing intrauterine devices
have been the most reliable contraceptive methods at population
scale for two decades. The design problem is unusual: the drug
release must be constant or near-constant for three to five years,
with a predictable approach to subtherapeutic levels. The release
kinetics are zero-order or near-zero-order (constant rate, set by
the rate-limiting polymer membrane), not Crank-style $\sqrt{t}$.
The engineering choice of membrane material and thickness is the
controlling design knob.

\section{Sterilisation and packaging}\pathtag{standard}
% Figure: sterilisation process compatibility matrix.
% Vol VI Ch 12 sec 12.5.
\begin{figure}[htbp]
\centering
\begin{tikzpicture}[
  cell/.style={draw=engblack, minimum width=2.0cm, minimum height=0.8cm,
    align=center, font=\scriptsize},
  hcell/.style={draw=engblack, minimum width=2.0cm, minimum height=0.8cm,
    align=center, font=\scriptsize\bfseries, fill=engblack!10},
  ok/.style={fill=englime!30},
  caution/.style={fill=engblue!20},
  bad/.style={fill=engred!25},
]

% Header row
\node[hcell] at (0, 0) {Material};
\node[hcell] at (2.0, 0) {Autoclave\\(121--134\,$^{\circ}$C steam)};
\node[hcell] at (4.0, 0) {Ethylene\\oxide (EtO)};
\node[hcell] at (6.0, 0) {Gamma\\(25--40\,kGy)};
\node[hcell] at (8.0, 0) {Electron\\beam};
\node[hcell] at (10.0, 0) {Hydrogen\\peroxide gas};

% Stainless steel
\node[hcell] at (0, -0.8) {Stainless steel};
\node[cell, ok] at (2.0, -0.8) {Compatible};
\node[cell, ok] at (4.0, -0.8) {Compatible};
\node[cell, ok] at (6.0, -0.8) {Compatible};
\node[cell, ok] at (8.0, -0.8) {Compatible};
\node[cell, ok] at (10.0, -0.8) {Compatible};

% Co-Cr-Mo
\node[hcell] at (0, -1.6) {Co-Cr-Mo alloy};
\node[cell, ok] at (2.0, -1.6) {Compatible};
\node[cell, ok] at (4.0, -1.6) {Compatible};
\node[cell, ok] at (6.0, -1.6) {Compatible};
\node[cell, ok] at (8.0, -1.6) {Compatible};
\node[cell, ok] at (10.0, -1.6) {Compatible};

% Ti-6Al-4V
\node[hcell] at (0, -2.4) {Ti-6Al-4V};
\node[cell, ok] at (2.0, -2.4) {Compatible};
\node[cell, ok] at (4.0, -2.4) {Compatible};
\node[cell, ok] at (6.0, -2.4) {Compatible};
\node[cell, ok] at (8.0, -2.4) {Compatible};
\node[cell, ok] at (10.0, -2.4) {Compatible};

% UHMWPE
\node[hcell] at (0, -3.2) {UHMWPE};
\node[cell, caution] at (2.0, -3.2) {Avoid (creep)};
\node[cell, ok] at (4.0, -3.2) {Compatible};
\node[cell, bad] at (6.0, -3.2) {Oxidation};
\node[cell, bad] at (8.0, -3.2) {Oxidation};
\node[cell, ok] at (10.0, -3.2) {Compatible};

% PMMA
\node[hcell] at (0, -4.0) {PMMA (bone cement)};
\node[cell, bad] at (2.0, -4.0) {Bad (softens)};
\node[cell, ok] at (4.0, -4.0) {Compatible};
\node[cell, caution] at (6.0, -4.0) {Discolour};
\node[cell, caution] at (8.0, -4.0) {Discolour};
\node[cell, ok] at (10.0, -4.0) {Compatible};

% PLA / PLGA
\node[hcell] at (0, -4.8) {PLA / PLGA};
\node[cell, bad] at (2.0, -4.8) {Hydrolysis};
\node[cell, ok] at (4.0, -4.8) {Compatible};
\node[cell, caution] at (6.0, -4.8) {Chain scission};
\node[cell, caution] at (8.0, -4.8) {Chain scission};
\node[cell, ok] at (10.0, -4.8) {Compatible};

% Silicone
\node[hcell] at (0, -5.6) {Silicone};
\node[cell, ok] at (2.0, -5.6) {Compatible};
\node[cell, caution] at (4.0, -5.6) {EtO residual};
\node[cell, ok] at (6.0, -5.6) {Compatible};
\node[cell, ok] at (8.0, -5.6) {Compatible};
\node[cell, ok] at (10.0, -5.6) {Compatible};

% PEEK
\node[hcell] at (0, -6.4) {PEEK};
\node[cell, ok] at (2.0, -6.4) {Compatible};
\node[cell, ok] at (4.0, -6.4) {Compatible};
\node[cell, ok] at (6.0, -6.4) {Compatible};
\node[cell, ok] at (8.0, -6.4) {Compatible};
\node[cell, ok] at (10.0, -6.4) {Compatible};

% Electronics / battery
\node[hcell] at (0, -7.2) {Electronics, battery};
\node[cell, bad] at (2.0, -7.2) {Bad (thermal)};
\node[cell, ok] at (4.0, -7.2) {Compatible};
\node[cell, caution] at (6.0, -7.2) {Dose limited};
\node[cell, caution] at (8.0, -7.2) {Dose limited};
\node[cell, ok] at (10.0, -7.2) {Compatible};

\end{tikzpicture}
\caption[Sterilisation-process compatibility matrix]{A first-cut
compatibility matrix between common device materials and the five
production-scale terminal sterilisation processes used in
medical-device manufacture, current as of 2026. Green cells
indicate routine compatibility under standard process parameters;
blue cells indicate compatibility with specific cautions
(monitoring, sub-saturating dose, residual outgassing); red cells
indicate the process is not appropriate without protective
measures (encapsulation, lower-temperature variant, or alternative
method). The matrix does not capture cycle-specific parameters
(temperature, dwell time, dose, humidity) and is a first-cut
selector, not a process specification. Standard references for
the specific processes are \cite{std:v6c12-iso-11135-2014} (EtO),
\cite{std:v6c12-iso-11137-1-2006} (radiation), and
\cite{text:v6c12-ratner-biomaterials-2020} for the material-
specific responses.}
\label{fig:vol06:ch12:sterilisation-matrix}
\end{figure}


A medical device that contacts a patient or a sterile body cavity
must be sterile at the point of use. The international consensus
on sterility, as a design and regulatory target, is the
\engterm{sterility assurance level} (SAL) of $10^{-6}$: at most
one in a million devices in a sterilised lot carries a viable
microorganism. The target is achieved by terminal sterilisation
(after final packaging) or by aseptic processing (combining
sterile components in a sterile environment), depending on the
device's material and assembly constraints.

\subsection{Terminal sterilisation processes}

\engterm{Autoclave (saturated steam)} sterilises by heat and
moisture. Standard cycles run at 121 $^{\circ}$C for 15 to 30
minutes or at 134 $^{\circ}$C for 3 to 18 minutes, depending on
the load configuration. Autoclave is fast, low-cost, residue-free,
and demands a material that survives saturated steam above 121
$^{\circ}$C. Metallic instruments (stainless, cobalt-chromium,
titanium) tolerate autoclaving routinely. Most polymeric devices
and all electronic devices use a different terminal process,
because UHMWPE creeps, PMMA softens, PLA hydrolyses, and
polyolefins may deform at autoclave temperatures.

\engterm{Ethylene oxide (EtO)} sterilises by alkylation of microbial
DNA. A standard cycle exposes the device at 30 to 60 $^{\circ}$C
to an EtO concentration of 600 to 1200 mg/L for 2 to 6 hours, with
controlled humidity (40 to 80 percent relative humidity to permit
EtO diffusion through the device's water-permeable barriers). The
cycle is followed by an aeration step (one to several days) to
remove residual EtO and its principal degradation products, EtO
chlorohydrin and ethylene glycol, to levels below ISO 10993-7
limits. EtO is compatible with most polymeric materials, with
silicone (with extended aeration), with electronics, and with
heat-sensitive biological components. The process is regulated by
ISO 11135 \cite{std:v6c12-iso-11135-2014}. The principal hazards
are EtO's flammability, its acute toxicity to operators, and the
need to validate the residuals against the device's contact
duration and contact type.

\engterm{Gamma sterilisation} delivers a high-energy photon dose
(25 to 40 kGy is the standard sterilisation range; the dose is set
by ISO 11137-1 for the bioburden of the device's manufacturing
environment) from a cobalt-60 source. Gamma is fast, penetrating
(useful for multi-layer packaging), and residue-free. The hazard is
material damage: ionising radiation produces free radicals that
oxidise UHMWPE, scission PLA and PMMA chains, discolour polymers,
and degrade some biologics. Gamma-sterilised UHMWPE in air was the
standard for decades and is now the canonical example of an
in-service degradation problem caused by the sterilisation
process; modern UHMWPE is gamma-sterilised in inert atmosphere and
thermally cross-linked to suppress oxidation.

\engterm{Electron-beam sterilisation} delivers the same standard
dose range as gamma (25 to 40 kGy) using a high-energy electron
accelerator. The beam is less penetrating than gamma (depth on the
order of a centimetre at standard energies) and is therefore
suited to thin, well-defined geometries. The material-compatibility
profile is similar to gamma.

\engterm{Hydrogen peroxide gas plasma} (the STERRAD process and
its descendants) sterilises by oxidative attack of microbial
molecules. The process runs at 50 to 60 $^{\circ}$C, compatible
with electronics and most polymers, with residual peroxide that
decays to water and oxygen. The process is the modern choice for
re-sterilisable hospital instruments that include electronics and
heat-sensitive polymers.

\engterm{Vapour-phase hydrogen peroxide (VHP)} and
\engterm{paracetic acid} fill specialised niches: re-processing of
sealed-source devices, terminal sterilisation of complex polymer-
electronic assemblies, sterilisation of clean rooms.

The sterilisation choice is not independent of the design. A
device that combines UHMWPE, electronics, and a degradable polymer
has no single terminal-sterilisation cycle compatible with all
three components; the design demands aseptic assembly of separately
sterilised modules, or a re-design that eliminates one of the
incompatible materials.

\subsection{Packaging as part of the sterile barrier system}

The packaging is the sterile barrier system. From the moment of
terminal sterilisation, the package is the only thing standing
between the device and the environment. ISO 11607 specifies the
sterile barrier system's design requirements:

\begin{itemize}[leftmargin=*]
\item Material that survives the sterilisation cycle without
  compromising the package integrity (Tyvek substrate for EtO and
  gamma; specific polymer films for autoclave; pouch construction
  for electron beam).
\item Seal integrity that resists pinholes and channel leakage
  under transport and storage.
\item Tamper-evident features.
\item Microbial-barrier integrity tested by bubble emission,
  dye penetration, or sealed-package microbial-challenge testing.
\item Shelf-life testing under accelerated and real-time
  conditions.
\item Aseptic-opening features that prevent contamination at the
  point of use.
\end{itemize}

A failure of the sterile barrier is a failure of the sterilisation:
a device that was sterile at the end of the cycle but contaminates
through a compromised package is not a sterile device at the point
of use. Sterile-barrier defect rates are tracked as a quality
metric.

\subsection{Bioburden and parametric release}

The sterilisation cycle is validated against the device's
manufacturing bioburden. ISO 11737-1 specifies the bioburden test:
the count of viable microorganisms on a sample of devices taken at
the end of the manufacturing line. The bioburden number sets the
required logarithmic reduction. A device manufactured with average
bioburden of 100 colony-forming units per device, sterilised to
SAL $10^{-6}$, requires an eight-log reduction. The validated
sterilisation cycle delivers that reduction with a safety margin.

\engterm{Parametric release} is the regulatory option to release
sterile-product lots based on the validated process parameters
(temperature, time, dose, gas concentration) rather than on
end-product sterility testing. The option requires demonstrated
process stability and is more common for high-volume EtO and
gamma cycles than for steam or low-temperature gas cycles.

\section{Regulatory framework (FDA, EMA) for medical devices}\pathtag{core}
% Figure: FDA medical-device clearance pathway flowchart.
% Vol VI Ch 12 sec 12.6.
\begin{figure}[htbp]
\centering
\begin{tikzpicture}[
  start/.style={draw=engblack, fill=engblue!15, minimum height=0.9cm,
    minimum width=3.5cm, align=center, font=\small\bfseries},
  decision/.style={draw=engblack, fill=englime!10, diamond,
    aspect=2.2, align=center, font=\scriptsize},
  pathway/.style={draw=engblack, fill=engblue!10, minimum height=1.0cm,
    minimum width=3.0cm, align=center, font=\small\bfseries},
  outcome/.style={draw=engblack, fill=engblack!10, minimum height=0.7cm,
    minimum width=3.0cm, align=center, font=\scriptsize},
  arrow/.style={->, thick, engblack},
  label/.style={font=\scriptsize, anchor=west},
]

% Start
\node[start] (start) at (0, 7.5) {Device intended for human use};

% First decision: classification
\node[decision] (class) at (0, 6.0) {Risk class\\(Class I/II/III)};

% Class I path
\node[pathway, fill=englime!10] (classI) at (-6, 4.5) {Class I\\(general controls)};
\node[outcome] (i-out) at (-6, 3.2) {Most exempt\\from 510(k)};

% Class II path with 510(k)
\node[pathway, fill=engblue!10] (classII) at (-2, 4.5) {Class II\\(special controls)};
\node[decision] (predicate) at (-2, 3.0) {Predicate\\available?};
\node[pathway, fill=engblue!15] (k510) at (-3.5, 1.5) {510(k)\\substantial equivalence};
\node[pathway, fill=engred!10] (denovo) at (-0.5, 1.5) {De Novo\\classification};
\node[outcome] (k-out) at (-3.5, 0.2) {Cleared\\(~3,500/yr 2023)};
\node[outcome] (d-out) at (-0.5, 0.2) {Granted};

% Class III path with PMA
\node[pathway, fill=engred!15] (classIII) at (3, 4.5) {Class III\\(highest risk)};
\node[pathway, fill=engred!20] (pma) at (3, 3.0) {PMA\\(premarket approval)};
\node[outcome] (pma-out) at (3, 1.7) {Original PMAs:\\~20--40/yr};
\node[outcome] (pma-sup) at (3, 0.5) {PMA supplements:\\hundreds/yr};

% Post-market
\node[pathway, fill=engblack!10, minimum width=11.5cm, minimum height=0.7cm]
  (postmkt) at (-1.5, -1.2) {Post-market surveillance:
  MAUDE adverse-event reporting, Section 522 studies,
  recall classes I/II/III};

% Arrows
\draw[arrow] (start) -- (class);
\draw[arrow] (class) -- node[label, font=\scriptsize] {low} (classI);
\draw[arrow] (class) -- node[label, font=\scriptsize] {moderate} (classII);
\draw[arrow] (class) -- node[label, font=\scriptsize] {high} (classIII);

\draw[arrow] (classI) -- (i-out);

\draw[arrow] (classII) -- (predicate);
\draw[arrow] (predicate) -- node[label, font=\scriptsize] {yes} (k510);
\draw[arrow] (predicate) -- node[label, font=\scriptsize] {no} (denovo);
\draw[arrow] (k510) -- (k-out);
\draw[arrow] (denovo) -- (d-out);

\draw[arrow] (classIII) -- (pma);
\draw[arrow] (pma) -- (pma-out);
\draw[arrow] (pma-out) -- (pma-sup);

\draw[arrow] (i-out) |- (postmkt);
\draw[arrow] (k-out) |- (postmkt);
\draw[arrow] (d-out) |- (postmkt);
\draw[arrow] (pma-sup) |- (postmkt);

\end{tikzpicture}
\caption[FDA medical-device clearance pathways]{The three main
FDA pathways to market for a medical device, schematic. Class I
devices (bandages, examination gloves, manual stethoscopes) are
subject to general controls and are largely exempt from 510(k)
review. Class II devices (most infusion pumps, surgical needles,
many imaging accessories) take the 510(k) substantial-equivalence
pathway if a predicate device exists, or the De Novo
classification pathway if the device is novel-low-to-moderate
risk. Class III devices (life-sustaining or life-supporting
implants such as mechanical heart valves and pacemakers) require
the Premarket Approval (PMA) pathway, which requires clinical-
outcome data. Counts are current as of 2023 from the FDA 510(k)
and PMA databases. Post-market surveillance applies to all
classes through MAUDE adverse-event reporting, Section 522 post-
market studies, and the recall taxonomy I (most serious) /
II / III \cite{law:v6c12-21cfr-807, law:v6c12-21cfr-814,
law:v6c12-21usc-360c, web:v6c12-fda-510k-database-2024,
web:v6c12-fda-pma-database-2024}.}
\label{fig:vol06:ch12:fda-pathway}
\end{figure}


The medical-device regulatory framework distinguishes itself from
the drug regulatory framework in two principal ways: it is younger
(in the United States, dating from 1976; in Europe, principally
from 1993), and it is structured around risk classes rather than
around per-product clinical trials. The framework reflects a
practical compromise: clinical trial of every new device at full
drug-trial depth would be prohibitively slow and expensive given
the iterative nature of device improvement; absence of any
clinical trial would surrender the public to unmonitored failure.
The risk-class framework attempts to allocate clinical-evidence
burden to the devices that pose the greatest risk.

\subsection{United States: the FDA framework}

The United States Food and Drug Administration (FDA) regulates
medical devices under the 1938 Federal Food, Drug, and Cosmetic
Act, as amended by the 1976 Medical Device Amendments
\cite{law:v6c12-fdca-1938}. The 1976 amendments created three
classes for medical devices:

\begin{itemize}[leftmargin=*]
\item \engterm{Class I} (low risk): devices subject to general
  controls only (registration, listing, labelling, good
  manufacturing practice). Examples: bandages, examination gloves,
  manual stethoscopes. Most Class I devices are exempt from
  premarket review.
\item \engterm{Class II} (moderate risk): devices subject to
  general controls plus special controls (performance standards,
  post-market surveillance, special labelling, premarket
  notification). Examples: infusion pumps, surgical needles,
  diagnostic ultrasound transducers. Most Class II devices require
  premarket notification under section 510(k) of the FDCA.
\item \engterm{Class III} (high risk): devices subject to general
  controls plus Premarket Approval (PMA). Devices that support or
  sustain human life, are of substantial importance in preventing
  impairment of human health, or present a potential unreasonable
  risk of illness or injury. Examples: mechanical heart valves,
  pacemakers, implanted neurostimulators, deep-brain stimulators.
\end{itemize}

The class is assigned by 21 USC 360c on the basis of the device's
generic type, not on the basis of the specific product
\cite{law:v6c12-21usc-360c}.

\subsection{The 510(k) substantial-equivalence pathway}

The 510(k) pathway, codified in 21 CFR Part 807 Subpart E
\cite{law:v6c12-21cfr-807}, allows a manufacturer to market a
device by demonstrating substantial equivalence to a
\engterm{predicate device} already on the market. The substantial-
equivalence determination requires that the new device have the
same intended use as the predicate and either the same
technological characteristics or different characteristics that do
not raise different safety and effectiveness questions. The
manufacturer submits comparative data (often bench testing,
sometimes animal testing, infrequently clinical-outcome data) and
the FDA reviews against a standard timeline (the FDA's
performance goal is 90 days for review). In fiscal year 2023, the
FDA cleared approximately 3,500 devices via 510(k) (current as of
2024) \cite{web:v6c12-fda-510k-database-2024}.

The 510(k) pathway has been the workhorse of United States
medical-device approval since 1976 and is also one of the most
critiqued elements of the framework. The Institute of Medicine's
2011 report on the 510(k) recommended that the pathway be
replaced \cite{paper:v6c12-iom-510k-2011}; the IOM's critique was
that the pathway's predicate-device chain accumulates over
decades, allowing devices to come to market through a chain of
substantial-equivalence determinations to ancestors whose own
clinical record is fragmentary. The DePuy ASR case exemplifies
the failure mode: the ASR XL was cleared via 510(k) in 2005 with
a predicate chain extending back to earlier MoM components whose
own long-term outcome record did not support the substantial-
equivalence determination as applied to the small-clearance
edge-loading geometry that ultimately produced ASR's distinctive
failure mode \cite{paper:v6c12-cohen-mom-bmj-2012,
paper:v6c12-zaaijer-510k-2013,
paper:v6c12-makary-recalls-2017}.

\subsection{The PMA pathway}

The Premarket Approval pathway, codified in 21 CFR Part 814
\cite{law:v6c12-21cfr-814}, requires the manufacturer to
demonstrate reasonable assurance of safety and effectiveness
through clinical-outcome data, typically from a controlled
clinical investigation. PMA reviews are substantially longer than
510(k) reviews (medians on the order of 6 to 18 months for the
agency portion, plus 1 to 3 years for the clinical trial).
Original PMA approvals number in the tens per year (current as of
2024) \cite{web:v6c12-fda-pma-database-2024}; PMA supplements
(modifications to approved PMA devices) number in the hundreds per
year and reflect the iterative-improvement reality of the device
field.

\subsection{The De Novo pathway}

The De Novo classification pathway, established in 1997 and
substantially revised in 2012, applies to novel devices of
low-to-moderate risk for which no predicate exists. The De Novo
determination assigns the device to Class I or Class II and
provides a predicate for subsequent 510(k) submissions of similar
devices. De Novo numbers in the tens per year and is the principal
pathway by which a genuinely new device class enters the market
without requiring full PMA.

\subsection{Quality system and design controls}

Independently of the premarket pathway, every medical-device
manufacturer that places product in United States commerce must
comply with 21 CFR Part 820, the Quality System Regulation
\cite{law:v6c12-21cfr-820}. The QSR specifies requirements for:

\begin{itemize}[leftmargin=*]
\item Design controls (design input, design output, design
  review, design verification, design validation, design
  transfer, design change control).
\item Document controls.
\item Purchasing controls (supplier qualification).
\item Production and process controls.
\item Identification and traceability.
\item Inspection, measurement, and testing.
\item Corrective and preventive action (CAPA).
\item Complaint handling and medical-device reporting.
\end{itemize}

ISO 13485 \cite{std:v6c12-iso-13485-2016} is the international
counterpart to 21 CFR 820 and is the basis on which most
international notified bodies issue conformity certificates. The
two standards are sufficiently aligned that an ISO 13485-
compliant manufacturer typically meets most of the QSR with
modest gap-closure.

\subsection{Risk management and software life cycle}

ISO 14971 \cite{std:v6c12-iso-14971-2019} specifies the risk-
management process for medical devices: hazard identification,
risk analysis, risk evaluation, risk control, evaluation of
overall residual risk, risk-management report, and post-production
information loop. The 2019 edition closed earlier gaps around
benefit-risk analysis and post-market integration. IEC 62304
\cite{std:v6c12-iec-62304-2015} specifies the software life-
cycle process for medical-device software, with safety classes A
(non-injury), B (non-serious injury), and C (death or serious
injury). The classes drive proportionate development rigour. The
Therac-25 case (section 12.7) is the founding episode of medical-
device software discipline; IEC 62304's requirements would have
mandated formal hazard analysis, design review, and verification
testing of the safety-critical subsystem in which the race
condition lived.

\subsection{European Union: the MDR}

The European Union regulates medical devices under Regulation
(EU) 2017/745, the Medical Devices Regulation (MDR), which
replaced the 1993 Medical Devices Directive (MDD) effective May
2021 \cite{law:v6c12-eu-mdr-2017-745}. The MDR retains a risk-
class structure (Class I, IIa, IIb, III, with implantable devices
generally in IIb or III and active implantable devices in III)
and a Notified Body conformity-assessment route. The MDR
strengthens requirements for clinical evidence (particularly for
Class III and implantable devices), establishes the European
database on medical devices (EUDAMED, see
\cite{web:v6c12-eudamed-2024}), introduces Unique Device
Identification (UDI), and imposes post-market surveillance and
vigilance obligations.

The MDR's clinical-evidence requirements are tighter than the
510(k) pathway's. A Class III implantable device under MDR
generally requires a clinical-investigation report (under ISO
14155) or equivalent literature-based evidence; equivalence to a
prior device is allowed only under restrictive conditions and is
not the dominant route to market. The MDR also gives Notified
Bodies more inspection authority and gives competent authorities
in member states broader recall and market-withdrawal powers.

\subsection{Post-market surveillance and recalls}

Both the FDA and the EU framework rely on post-market surveillance
to catch failure modes that did not appear in premarket testing.
The principal United States post-market mechanism is MAUDE
(Manufacturer and User Facility Device Experience), the FDA's
adverse-event reporting database, established under the Safe
Medical Devices Act of 1990 \cite{law:v6c12-safe-medical-devices-act-1990}
and accessible at \cite{web:v6c12-fda-maude-2024}. MAUDE collects
reports from manufacturers, distributors, hospitals, and user
facilities under mandatory reporting timelines. The database is
public.

Recalls are classified into three classes by health-consequence
severity:

\begin{itemize}[leftmargin=*]
\item Class I recall: reasonable probability that use will cause
  serious adverse health consequences or death.
\item Class II recall: use may cause temporary or medically
  reversible adverse health consequences, or remote probability of
  serious adverse health consequences.
\item Class III recall: use is not likely to cause adverse health
  consequences.
\end{itemize}

The recall class is decided by the FDA in consultation with the
manufacturer. Class I recalls are the most serious and are
typically the focus of root-cause analyses by the manufacturer
and the regulator. Recent decadal trends in recall counts have
been upward, principally for Class II (the bulk of recalls), with
device-software-related recalls becoming a growing share
\cite{paper:v6c12-amato-recall-data-2014}. The synthesised
distribution in \texttt{data/fda\_recalls\_by\_class.csv} is
representative of the shape (Class II dominant, Class I a small
share, Class III rare), reproduced for the exercises in this
chapter.

The EU vigilance system parallels MAUDE. EUDAMED includes a
vigilance module for adverse-event reporting and a market-
surveillance module for competent-authority actions. The MHRA's
2010 and 2012 medical-device alerts on metal-on-metal hip implants
\cite{web:v6c12-mhra-mom-2012} exemplify the system in action.

\subsection{Joint registries as post-market evidence}

In the orthopaedic field, national joint registries have been the
single most informative post-market surveillance instrument. The
Swedish, Norwegian, Australian (AOANJRR), and UK NJR registries
have accumulated cohort sizes in the hundreds of thousands of
implants with annual revision-rate reporting by device, by
fixation method, and by patient demographic. The DePuy ASR signal
appeared in the AOANJRR 2007 report and the NJR 2010 report well
before the FDA recall \cite{web:v6c12-aoanjrr-2023,
web:v6c12-njr-2023, paper:v6c12-smith-asr-revision-2012}. The
registries are the principal reason the orthopaedic field
identifies failing devices on a one- to three-year horizon rather
than a decade-plus horizon
\cite{paper:v6c12-kessler-fdca-history-1996}.

\section{Failure: Therac-25 (software), Bjork-Shiley heart valve, hip implants}\pathtag{core}
% Figure: Bjork-Shiley convexo-concave valve strut fracture schematic.
% Vol VI Ch 12 sec 12.7.
\begin{figure}[htbp]
\centering
\begin{tikzpicture}[
  arrow/.style={->, thick, engblack},
  label/.style={font=\scriptsize, align=center},
]

% Valve ring viewed from outflow side
\node[font=\small\bfseries] at (-3.5, 4.5) {(a) Valve ring (outflow view)};
\draw[engblack, ultra thick] (-3.5, 2.0) circle (1.8);
\draw[engblue!30, very thick] (-3.5, 2.0) circle (1.7);

% Inlet strut (large) at bottom
\draw[engblack, ultra thick, fill=engblack!30]
  (-3.5, 2.0) to[bend left=20] (-3.5, 0.5)
  to[bend right=20] (-3.5, 2.0);
\node[label] at (-3.5, 0.0) {Inlet strut\\(integral)};

% Outlet strut at top: the failure-prone element
\draw[engred, ultra thick, fill=engred!30]
  (-3.5, 2.0) to[bend right=20] (-3.5, 3.5)
  to[bend left=20] (-3.5, 2.0);
\node[label, engred] at (-3.5, 4.1) {Outlet strut\\(welded)};

% Weld marks on outlet strut
\fill[engred!70] (-3.55, 3.4) circle (0.04);
\fill[engred!70] (-3.45, 3.4) circle (0.04);

% Disc indicator
\draw[engblue, dashed] (-3.5, 1.6) circle (1.55);
\node[label, engblue] at (-3.5, 1.6) {Pyrolytic-\\carbon disc};

% Detail (b): close-up of weld fatigue at the strut foot
\node[font=\small\bfseries] at (3.5, 4.5) {(b) Weld-toe fatigue (close-up)};

% Ring base
\draw[engblack, ultra thick, fill=engblack!20]
  (1.5, 1.5) -- (5.5, 1.5) -- (5.5, 2.0) -- (1.5, 2.0) -- cycle;
\node[label] at (3.5, 1.2) {Haynes-25 ring (cross-section)};

% Strut foot
\draw[engblack, ultra thick, fill=engblack!30]
  (3.2, 2.0) -- (3.8, 2.0) -- (3.8, 3.8) -- (3.2, 3.8) -- cycle;
\node[label] at (3.5, 4.0) {Strut leg};

% Weld bead
\fill[engred!50] (3.0, 1.95) .. controls (3.0, 2.3) and (3.2, 2.4) ..
  (3.4, 2.4) -- (3.6, 2.4) ..
  controls (3.8, 2.4) and (4.0, 2.3) .. (4.0, 1.95) -- cycle;
\node[label, engred, anchor=west] at (4.1, 2.2) {Weld bead};

% Crack initiation site at weld toe
\draw[engred, very thick] (3.2, 2.3) -- (3.1, 2.5) -- (3.05, 2.6);
\draw[engred, very thick] (3.8, 2.3) -- (3.9, 2.5);

% Annotation
\draw[arrow, engred] (5.5, 2.7) -- (3.2, 2.5);
\node[label, engred, anchor=west, align=left] at (5.5, 3.0)
  {Fatigue crack initiates\\at weld toe\\(stress concentration)};

% Cycle counter annotation
\node[label, anchor=west, align=left] at (5.5, 0.5)
  {$\approx 35\text{--}40 \times 10^{6}$ cycles\\per patient-year\\
  ($\approx 70$ bpm)};

\end{tikzpicture}
\caption[Bjork-Shiley convexo-concave outlet strut and weld-toe
fatigue]{(a) Outflow-side view of the Bjork-Shiley convexo-
concave (BSCC) valve ring, showing the integral inlet strut
(bottom, black) and the welded outlet strut (top, red). The
pyrolytic-carbon occluder disc (dashed blue) opens to 60 or 70
degrees depending on the design variant. (b) Cross-sectional
detail of the outlet-strut weld foot. The weld bead (red) joins
the strut leg to the Haynes-25 ring; the geometric discontinuity
at the weld toe is a stress concentrator. Under the cyclic
bending load applied by the disc on every cardiac cycle
(approximately 35 to 40 million cycles per patient-year at 70
beats per minute), fatigue cracks initiate at the weld toe and
propagate to single-leg failure (SLF), then to complete strut
fracture and disc embolisation. The schematic is based on the
metallurgical reconstruction reported by van der Lugt and
others (1993) and the long-term cohort analysis of Blot,
Ibrahim, Inglesby and others (2005)
\cite{paper:v6c12-vanderlugt-strut-fracture-1993,
paper:v6c12-blot-bjork-shiley-2005,
paper:v6c12-paris-erdogan-1963}.}
\label{fig:vol06:ch12:bjork-shiley-strut}
\end{figure}


Three named cases organise this chapter's failure section. Each
illustrates a different failure mode, a different organisational
mechanism, and a different regulatory response. Together they
sketch the modern field of medical-device safety as a discipline.

\subsection{Therac-25 (1985-1987): the migration of safety from hardware to software}

Therac-25, a radiation-therapy machine manufactured by Atomic
Energy of Canada Limited and sold from 1982, delivered massive
radiation overdoses to six patients between June 1985 and January
1987 \cite{acc:therac-25-1985-1987, paper:leveson-turner1993}.
Three of the six events were fatal. The closest-equivalent
canonical source for the case is the 1993 IEEE Computer paper of
Leveson and Turner, supplemented by Leveson's 2011
\emph{Engineering a Safer World}
\cite{text:v6c12-leveson-engineering-safer-world-2011}.

\engterm{The technical mechanism}. Therac-25 supported two beam
modes: a low-current electron mode for surface treatment, and a
high-current X-ray mode in which a metal scattering target and
tungsten flattener intervened between a high-energy electron beam
and the patient. Switching between modes required mechanical
movement of the target into or out of the beam path. The control
software contained race conditions that allowed the operator
console to indicate the X-ray mode while the mechanical target had
not moved into position, or allowed the high-energy electron beam
to fire when the operator believed the low-energy electron mode
was selected. In those conditions the patient received the
high-energy electron beam without the target's scattering, a
factor of approximately $10^2$ higher than the prescribed dose at
the entry surface \cite{paper:leveson-turner1993}.

\engterm{The organisational mechanism}. The race condition was not
detected by the machine's safety interlocks. The interlocks had
been migrated from earlier Therac models (Therac-6 and Therac-20)
that had implemented hardware interlocks for the same hazard.
Therac-25 replaced the hardware interlocks with software, on the
assumption that the underlying software had been validated by the
prior models' service history. The assumption was incorrect: the
earlier models' hardware backups had masked software errors that
became unmasked when Therac-25 removed the hardware. Operator
displays showed terse messages such as ``MALFUNCTION 54'' without a
safety interpretation. Operators in some sites, on the basis that
prior Therac machines' malfunction codes had typically indicated
recoverable beam-positioning errors, used a ``P'' command to
proceed past the malfunction. In Therac-25 the same display
masked a dose anomaly; proceeding through the malfunction allowed
the beam to fire again with the same misconfiguration. The Tyler
dose display, in at least one case, showed an underdose indication
while the patient had received a massive overdose; the operator
restarted the treatment without realising the patient had already
been irradiated. AECL's response chain was slow to recognise a
recurring fault; the same failure mode recurred across multiple
sites for nearly two years before the machine was withdrawn in
January 1987.

\engterm{The regulatory response}. No single official accident-
investigation report exists. The FDA acted as a regulator monitoring
AECL's response and required design changes following each event;
the Canadian Health Protection Branch acted similarly. The
canonical reconstruction is Leveson and Turner's 1993 paper. The
case became a founding episode in safety-critical software
engineering and is one of the two or three most-cited cases in the
medical-device software literature. The 2006 IEC 62304 software
life-cycle standard (and its 2015 amendment, current edition,
\cite{std:v6c12-iec-62304-2015}) requires a hazard analysis that
recognises the removal of a hardware interlock as a safety-critical
change, an independent software review for safety-critical paths,
operator-display design that distinguishes routine malfunctions
from safety anomalies, and a vigilance feedback loop that catches
recurring faults.

\engterm{The engineering lesson}. The Therac-25 case is the
migration of safety from hardware to software without the hazard
analysis the migration required. The hardware interlocks of
Therac-6 and Therac-20 were load-bearing safety functions.
Removing them shifted the load to software that had not been
independently validated against the hazards the hardware had been
catching. The cost was paid by patients. The 1990 Safe Medical
Devices Act \cite{law:v6c12-safe-medical-devices-act-1990}
extended adverse-event reporting obligations to user facilities
and gave the FDA new authority over post-market surveillance
partly in response to the recognition that medical-device software
failures had been under-reported through the 1980s.

\subsection{Bjork-Shiley convexo-concave heart valve (1979-1986): fatigue, weld defects, and the disclosure failure}

The Bjork-Shiley convexo-concave (BSCC) mechanical heart valve,
sold by Shiley Inc. from 1979 until withdrawal from the United
States market in November 1986, fractured at the welded outlet
strut at a rate that depended on valve size, opening angle, and
implant orientation \cite{acc:bjork-shiley-cc-1979-1986,
paper:v6c12-blot-bjork-shiley-2005,
web:v6c12-fda-bjork-shiley-1992-recall}. Approximately 86,000
valves were implanted worldwide; the cumulative count of confirmed
strut fractures reached 619 by 2003 with roughly 400 deaths
\cite{paper:v6c12-blot-bjork-shiley-2005}.

\engterm{The technical mechanism}. The BSCC valve consists of a
pyrolytic-carbon occluder disc, a Haynes-25 (cobalt-chromium-
tungsten) metal ring, and two welded struts that retain the disc.
The convexo-concave geometry, introduced in 1979 as a hemodynamic
improvement on the earlier flat-disc Bjork-Shiley model, opened the
disc to a larger angle (60 to 70 degrees, depending on variant) to
reduce regurgitation. The outlet strut, the smaller of the two
retention legs, was welded to the ring at two locations. Under
the cyclic bending load applied by the disc on every cardiac
cycle (approximately 35 to 40 million cycles per patient-year),
fatigue cracks initiated at the weld toes and propagated until
the strut fractured in a single ligament. Once the second leg
failed, the disc embolised, the patient suffered acute massive
regurgitation, and mortality from a single-leg-then-strut-fracture
event was approximately 70 percent without immediate emergency
valve replacement \cite{paper:v6c12-blot-bjork-shiley-2005,
paper:v6c12-vanderlugt-strut-fracture-1993}. The fracture risk was
strongly size-dependent: the 29 mm and larger mitral valves with
70-degree opening showed annual strut-fracture incidence of 0.20
to 0.50 percent per patient-year, while smaller aortic valves
with 60-degree opening showed rates an order of magnitude lower.
The mechanism is captured by the Paris-Erdogan crack-growth law
applied at the weld-toe stress concentration; the integrator in
\texttt{code/bjork\_shiley\_fracture\_paris.py} reproduces the
qualitative size dependence.

\engterm{The organisational mechanism}. A subset of valves carried
documented weld defects. Internal Shiley quality-control records,
later released in civil-litigation discovery, identified
out-of-specification welds and showed that some valves with
documented rework had not been segregated from the production
stream. The defect rate inferred from these records was not
uniform across production lots, and specific manufacturing-date
cohorts showed elevated fracture rates consistent with concentrated
defects \cite{web:v6c12-doj-shiley-1994-settlement}. The
convexo-concave variant had been approved as a supplement to the
existing Bjork-Shiley PMA rather than as a new device requiring
fresh trials; the supplement pathway did not require new clinical-
outcome data. The first strut fractures were reported to FDA in
1980. Shiley withdrew the larger 70-degree valves in 1983 and the
smaller 60-degree valves in 1986, and did not formally recall the
implanted cohort. The recommendation for implanted patients
remained surveillance rather than prophylactic explantation,
because the per-event mortality of elective re-operation in older
patients with comorbidities was, for many cohorts, comparable to
or higher than the per-year fracture risk.

\engterm{The regulatory response}. The 1992 FDA consent agreement
established the Supplemental Surgical Mortality Risk (SSMR)
calculator and a patient-notification programme. The 1994 U.S.
Department of Justice civil settlement of \$10{,}750{,}000 over
the undisclosed weld-defect records
\cite{web:v6c12-doj-shiley-1994-settlement} was one of the
founding episodes of the modern post-market device-surveillance
regime. The 1990 Safe Medical Devices Act
\cite{law:v6c12-safe-medical-devices-act-1990}, partly motivated
by the BSCC experience, strengthened FDA authority over post-market
surveillance and required adverse-event reporting from hospitals
and nursing homes. The BSCC case also drove the recognition that a
PMA supplement is not always a low-burden pathway: the supplement
mechanism for a geometric variant of a Class III implantable device
can introduce a novel failure mode (in this case, a weld-toe
fatigue mode in the new outlet-strut geometry) that the original
PMA's record does not anticipate.

\engterm{The engineering lesson}. The BSCC case has three layered
lessons. The first is fatigue: a cyclically loaded weld in a
load-bearing safety-critical component must be qualified at the
load spectrum and cycle count the in-service environment imposes,
with a quality system that identifies and rejects out-of-
specification welds. The second is recall versus surveillance:
when the per-event mortality of explantation exceeds the per-year
in-service hazard, recall is not the right response and surveillance
plus patient notification becomes the working alternative; the
notification system itself becomes a regulated engineering object.
The third is disclosure: the regulator's ability to act depends on
the integrity of the manufacturer's quality-control disclosure;
the 1990 SMDA and the 1996 Quality System Regulation revised the
expectations around manufacturer disclosure in response to the
BSCC experience.

\subsection{DePuy ASR (2010): metal-on-metal, edge loading, and the 510(k) blind spot}

The DePuy ASR XL Acetabular System (a total hip replacement) and
the DePuy ASR Hip Resurfacing System were manufactured by DePuy
Orthopaedics (Johnson \& Johnson) and sold worldwide from 2003 to
2010. DePuy issued a worldwide voluntary recall on 26 August 2010
after joint-registry data showed five-year cumulative revision
rates of 12 to 13 percent, against an expected rate of one to two
percent for a successful total hip replacement at the same horizon
\cite{acc:depuy-asr-hip-2010,
web:v6c12-fda-depuy-asr-2010-recall, paper:v6c12-cohen-mom-bmj-2012,
paper:v6c12-langton-asr-2011, paper:v6c12-smith-asr-revision-2012}.
Approximately 93,000 ASR units had been implanted worldwide. The
canonical baseline for comparison is Charnley's low-friction
arthroplasty \cite{paper:v6c12-charnley-low-friction-1972}; modern
total hip replacements with non-MoM bearings achieve revision
rates well under the ASR's per Poss and others
\cite{paper:v6c12-poss-hip-survival-2018}, and the structural
valve-deterioration analogue for valve patients is summarised by
Grunkemeier and others
\cite{paper:v6c12-grunkemeier-bioprosthetic-svd-2017}.

\engterm{The technical mechanism}. The ASR was a metal-on-metal
bearing: both the femoral head and the acetabular cup were
manufactured from high-carbon cobalt-chromium-molybdenum alloy
(ASTM F1537 or equivalent, \cite{std:v6c12-astm-f1537-2020,
std:v6c12-astm-f75-2018}). The bearing was designed with low
diametral clearance (approximately 80 to 150 $\mu$m) and a low
cup-rim coverage angle of approximately 144 degrees, the latter
chosen to reduce the cup wall thickness and allow larger femoral
heads in patients with smaller anatomy. The design intent was that
the small clearance would, in normal gait, support a hydrodynamic-
lubrication regime with very low wear
\cite{paper:v6c12-langton-asr-2011}. In service the lubrication
regime depended sensitively on cup inclination, anteversion, and
bearing geometry. When the cup was implanted at a steeper
inclination (above approximately 50 degrees from the horizontal),
or when anteversion fell outside a narrow window, the contact patch
migrated toward the cup rim. Edge loading then occurred during
routine gait, the lubrication film thinned or broke down, and
metal-on-metal contact produced submicrometre cobalt and chromium
wear particles. The local tissue response included aseptic
lymphocytic vasculitis-associated lesions (ALVAL), pseudotumours,
bone resorption around the implant, and elevated whole-blood
cobalt and chromium ion concentrations (cobalt frequently above
7 $\mu$g/L, with case reports above 1000 $\mu$g/L in extreme
failure) \cite{paper:v6c12-cohen-mom-bmj-2012,
paper:v6c12-jacobs-metal-ions-2003}. The cobalt and chromium
levels in \texttt{data/mom\_hip\_ion\_levels.csv} reproduce the
trend from Langton and others (2011) and from Cohen (2012).
Titanium-implant background on the comparator alloys is given in
ASTM F136 \cite{std:v6c12-astm-f136-2013}.

\engterm{The organisational mechanism}. The ASR XL was cleared in
the United States via the 510(k) substantial-equivalence pathway
in 2005. The predicate device chain extended back to earlier MoM
components that had not been subject to long-term clinical-outcome
registry comparison. The 510(k) review did not require new
clinical-outcome data, did not require finite-element or
hydrodynamic-lubrication analysis of the small-clearance geometry,
and did not require a registry-based post-market study. The 510(k)
clearance pathway as it existed in 2005 treated the ASR as
substantially equivalent to its predicate; the failure mode that
emerged in service (edge-loading-driven metal-ion release) was not
predicted from the predicate device's record
\cite{paper:v6c12-cohen-mom-bmj-2012,
paper:v6c12-zaaijer-510k-2013}. Joint registries accumulated faster
than DePuy or the regulators recognised the signal. The AOANJRR
flagged elevated ASR revision rates from its 2007 report onward
\cite{web:v6c12-aoanjrr-2023}. The NJR reported the ASR's five-year
cumulative revision rate at approximately 12 to 13 percent in its
2010 report. The UK MHRA issued a medical-device alert in April 2010
recommending annual blood-metal-ion testing for symptomatic patients
with MoM implants and tightened the alert in 2012
\cite{web:v6c12-mhra-mom-2012}. The FDA issued a safety
communication on metal-on-metal hip implants in January 2013
\cite{web:v6c12-fda-mom-orthopaedic-2013} and required post-market
surveillance studies for the broader MoM device class.

\engterm{The regulatory response}. The DePuy ASR recall was the
largest medical-device recall in orthopaedic history by patient
count. Subsequent United States multi-district litigation resulted
in settlements totalling more than \$2.5 billion (initial 2013
settlement) plus subsequent extensions. In 2016 the FDA
reclassified all metal-on-metal total hip replacement systems from
the 510(k) pathway to the PMA pathway under 21 USC 360e(b)(1), a
formal acknowledgement that the 510(k) substantial-equivalence
pathway was insufficient for the device class. A plain-language
overview of the 510(k) pathway, suitable for clinicians and for
trainee engineers, is available at
\cite{web:v6c12-ahrq-pcori-510k-explained-2018}.

\engterm{The engineering lesson}. The DePuy ASR case is the
canonical demonstration of the 510(k) pathway's blind spot for
design features whose failure mode is not present in the predicate's
failure record. The lubrication-regime collapse under edge loading
was not a failure mode of the predicate devices; it was a failure
mode of the specific ASR geometry, sensitive to surgeon implant-
position variability, that the 510(k) review did not capture and
that DePuy's premarket testing did not stress. The case also
demonstrates the post-market role of joint registries: the
AOANJRR's 2007 signal and the NJR's 2010 report identified the
failing device on a timescale measured in years rather than the
decade-plus timescale a litigation-driven evidence base would have
required. The joint registries are the principal reason the
orthopaedic field has a working post-market surveillance system
at all; the BMJ 2012 review describes the registry signal as the
single most informative element in the MoM hip story
\cite{paper:v6c12-cohen-mom-bmj-2012,
paper:v6c12-smith-asr-revision-2012}.

\subsection{The three cases as a unit}

The three cases share a structure. In each, a design choice that
appeared locally reasonable (replace hardware interlocks with
software; introduce a hemodynamically improved valve geometry;
introduce a low-clearance bearing for larger head sizes) produced
a failure mode that was not anticipated by the design's verification
process. In each, the failure recurred across multiple sites or
multiple patients before recognition. In each, the regulator's
response combined immediate corrective action with structural
change to the regulatory framework: the 1990 SMDA after BSCC
and Therac-25, the 2013 FDA MoM safety communication and 2016
reclassification after ASR, the IEC 62304 standard after Therac-25.

The reader who learns one general lesson from the three cases is
that biocompatibility, in the system-property sense given in
section 12.1, is not a property the manufacturer can self-certify
once at premarket. It is a property the system as a whole maintains
across the full service life through quality control, post-market
surveillance, professional reporting, and regulatory feedback.
The failure in each of these three cases was a failure of that
system, with the device itself the proximate-cause object but the
system the deeper cause.

\section{Estimation}\pathtag{core}

Two estimation blocks anchor this chapter. The first is mechanical
and lives in the implant world; the second is physico-chemical and
lives in the drug-device world.

\begin{estimation}
\textbf{Question.} Estimate the lifetime stress cycles accumulated
at the bearing surface of a total knee replacement in an active
patient, and the implied design margin against fatigue failure.

\textbf{Setup.} An active patient takes approximately 10{,}000
steps per day. The implant service horizon is 20 years. Each step
produces one major load cycle (heel strike to toe-off) at the
knee bearing surface plus approximately three minor cycles
associated with single-leg stance balance, for a working spectrum
of four cycles per step.

\textbf{Cycle count.}
\[
N = 10{,}000 \,\text{step}/\text{day}
\times 4 \,\text{cycles}/\text{step}
\times 365 \,\text{day}/\text{yr}
\times 20 \,\text{yr}
\approx 2.9 \times 10^8 \,\text{cycles}.
\]

\textbf{Stress amplitude.} A representative bearing-surface contact
stress is 30 to 80 MPa peak in walking, with a sliding component
producing a comparable shear at the polymer-on-metal interface.

\textbf{Material allowable.} Crosslinked UHMWPE acetabular and
tibial inserts have S-N curves in the literature with endurance
limits at approximately $10^7$ cycles around 20 to 30 MPa shear
amplitude; fully-reversed stress in this range produces measurable
crack initiation at $10^7$ to $10^8$ cycles. CoCr-on-CoCr and
CoCr-on-ceramic show endurance limits an order of magnitude higher
in stress and with different damage mechanisms (corrosion-fatigue
in the metal, fracture in the ceramic).

\textbf{Margin.} The implied lifetime $2.9 \times 10^8$ cycles is
at or above the threshold of the polymer's S-N curve. Crosslinked
UHMWPE bearings achieve clinical service lives of 15 to 25 years
in part because the actual stress amplitude at the bearing is
below the polymer's endurance limit and in part because the wear
mechanism dominates over fatigue. A bearing that operated at the
peak walking stress every cycle would not survive 20 years; the
real spectrum has a sub-yield mean stress and a peak that occurs
at a small duty cycle.

\textbf{Caveat.} The estimation conflates fatigue (crack-driven
failure) with wear (volume-loss failure). In the modern total
joint replacement, wear is the dominant failure mode by far. The
fatigue estimate sets a lower bound on what the bearing must
survive without crack initiation; the wear estimate sets the
expected service life.
\end{estimation}

\begin{estimation}
\textbf{Question.} Estimate the surface-area-to-volume ratio that
controls drug-elution kinetics on a coated coronary stent, and
the implied diffusivity required to deliver the drug over the
target clinical window.

\textbf{Setup.} A typical coronary stent has a length of 18 mm,
a deployed inner diameter of 3.0 mm, and a strut pattern that
covers approximately 12 percent of the lumen surface area. The
polymer drug-reservoir coating is 10 $\mu$m thick, applied uniformly
to the stent strut surfaces. The target clinical window is 30 days
of antiproliferative drug release.

\textbf{Coated surface area.}
\[
A_{\text{stent}} = \pi D L = \pi \times 3.0\,\text{mm}
\times 18\,\text{mm} = 170 \,\text{mm}^2.
\]
\[
A_{\text{coated}} \approx 0.12 \times 170 \,\text{mm}^2
= 20 \,\text{mm}^2.
\]

\textbf{Coating volume.}
\[
V_{\text{coat}} \approx A_{\text{coated}} \times L_{\text{coat}}
= 20 \,\text{mm}^2 \times 10\,\mu\text{m}
= 2 \times 10^{-4}\,\text{cm}^3.
\]

\textbf{Surface-to-volume ratio.}
\[
A/V = 1/L_{\text{coat}} = 1/(10\,\mu\text{m})
= 10^3 \,\text{cm}^{-1}.
\]

The thin-film geometry means that diffusion is one-dimensional
through $L_{\text{coat}}$; the surface area appears only as a
multiplier on the total quantity released, not in the
characteristic time.

\textbf{Characteristic time.} For a Crank-style 1D slab
$\tau = L^2 / (\pi^2 D)$. To achieve $\tau \approx 30 \,\text{days}
= 2.6 \times 10^6 \,\text{s}$ with $L = 10 \,\mu\text{m} = 10^{-3}
\,\text{cm}$,
\[
D \approx L^2 / (\pi^2 \tau)
= (10^{-3})^2 / (\pi^2 \times 2.6 \times 10^6)
\approx 4 \times 10^{-14} \,\text{cm}^2/\text{s}.
\]

\textbf{Sanity check.} Published polymer-drug diffusivities for
sirolimus and everolimus in PEVA/PBMA matrices are in the range
$10^{-13}$ to $10^{-12}$ cm$^2$/s; the estimate is one to two
orders of magnitude lower than the upper end of this range. The
discrepancy is principally because the closed-form Crank solution
treats the drug-tissue boundary as a perfect sink, while in
practice the tissue concentration builds and the boundary is
imperfect; the effective characteristic time is longer than the
slab model predicts. A diffusivity of $10^{-13}$ cm$^2$/s gives
$\tau \approx 100$ days at $L = 10 \,\mu\text{m}$, consistent
with the in vivo elution profile of modern sirolimus stents.

\textbf{Caveat.} The estimate is order-of-magnitude. Real drug-
eluting stent design uses the closed-form model as a first-cut
sizing tool and then confirms with bench elution, ex vivo
perfusion, and in vivo pharmacokinetics, as discussed in section
12.4.
\end{estimation}

\section{Project}\pathtag{core}

\begin{project}[Read a real device 510(k) submission and identify the hidden engineering choices]
The reader performs a desk-research analysis of a real,
publicly-available device 510(k) submission.

\textbf{Track.} Hybrid: document analysis plus engineering
reconstruction.

\textbf{Hazard class.} Low. No patient contact, no laboratory
work; the project is reading.

\textbf{Source.} The FDA's 510(k) database
\cite{web:v6c12-fda-510k-database-2024} is searchable by device
type, manufacturer, and clearance year. Pick a Class II device for
which a 510(k) summary is publicly accessible (most are). Pick a
device that interests you: an infusion pump, a continuous
glucose monitor, a sleep apnoea device, a surgical instrument,
a diagnostic imaging accessory.

\textbf{Specification.}

\begin{enumerate}[leftmargin=*]
\item \textbf{Identify the device.} 510(k) number, manufacturer,
  device name, product code, intended use, clearance year.
\item \textbf{Identify the predicate.} The submission states the
  predicate device(s). Find the predicate's own 510(k) record.
  Trace the predicate chain back two or three generations. Note
  the year of the earliest predicate in the chain.
\item \textbf{Read the substantial-equivalence statement.} What
  does the submission claim is the same as the predicate? What
  does it claim is different? For each difference, what testing
  was offered to support the claim that the difference does not
  raise different safety or effectiveness questions?
\item \textbf{Identify hidden engineering choices.} For each
  difference between the new device and the predicate, name the
  engineering choice the manufacturer made and the alternative the
  manufacturer rejected. Examples: changed sensor accuracy
  specification (what was the prior accuracy, what is the new
  accuracy, what testing supported the new accuracy?); changed
  battery chemistry (what was the prior chemistry, what is the
  new chemistry, what does the new chemistry imply for service life,
  end-of-life behaviour, recall logistics?); changed software
  architecture (what was the prior architecture, what is the new
  architecture, what is the IEC 62304 safety class?).
\item \textbf{Identify the biological-evaluation programme.} Which
  ISO 10993 tests were performed? Which were leveraged from the
  predicate? Which were waived? On what justification?
\item \textbf{Identify the sterilisation choice.} Process, cycle
  parameters, sterility assurance level, packaging.
\item \textbf{Read the labelling.} The 510(k) summary often includes
  the device's labelling. What contraindications, warnings, and
  precautions does the labelling assert? Which of these encode
  failure modes the manufacturer recognised but did not engineer
  out?
\item \textbf{Search MAUDE for the device.} The FDA MAUDE database
  \cite{web:v6c12-fda-maude-2024} carries adverse-event reports by
  device, by manufacturer, and by product code. How many events
  have been reported for this device since clearance? What are the
  three most common event types?
\item \textbf{Write the forensic note.} In two to three pages,
  reconstruct: (a) what the manufacturer claims the device does;
  (b) what the predicate chain suggests the device family has
  done; (c) what the post-market record says about the device's
  actual performance; (d) the three hidden engineering choices
  you identified, with their consequences; (e) the one
  recommendation you would make to the FDA reviewer if you were
  reviewing this submission today.
\end{enumerate}

\textbf{Deliverable.} A 6- to 10-page engineering notebook report
with the 510(k) number, the predicate chain, the hidden
engineering choices, the MAUDE summary, and the reviewer
recommendation. Cite specific 510(k) summary pages and MAUDE event
records.

\textbf{Discussion.} The point of the project is not to find a
failing device. Most 510(k)-cleared devices perform well. The
point is to read the regulatory record critically and to recognise
that every 510(k) submission embeds engineering decisions that the
substantial-equivalence framing tends to hide. The reader who
completes the project gains a working sense of how medical-device
design and medical-device regulation interact in practice.
\end{project}

\section{Exercises}

\subsection*{Calculation}\pathtag{core}

\begin{exercise}
A long-term implantable device contacts subcutaneous tissue
continuously for ten years. Using the ISO 10993-1 contact-type
and contact-duration categorisation, list the biological-evaluation
tests required for the device's premarket biological-evaluation
report.
\end{exercise}

\paragraph{Solution sketch.}
Contact type: implant device (in tissue). Contact duration: long-term
($> 30$ days). The ISO 10993-1 matrix for implant/long-term
requires cytotoxicity, sensitisation, irritation or intracutaneous
reactivity, acute systemic toxicity, subacute and chronic toxicity,
genotoxicity, implantation (the central test for implant devices),
and carcinogenicity. If the device contacts blood (vascular access,
blood-contacting implant), hemocompatibility is added. See
\cite{std:v6c12-iso-10993-1-2018} and the matrix reproduced in
\texttt{data/iso\_10993\_test\_matrix.csv}.

\begin{exercise}
A coronary stent strut is coated with a polymer reservoir
containing $50 \,\mu\text{g}$ of an antiproliferative drug, applied
to a coating of thickness $L = 8 \,\mu\text{m}$. Polymer-drug
diffusivity is $D = 2 \times 10^{-13} \,\text{cm}^2/\text{s}$.
Compute the characteristic release time $\tau = L^2 / (\pi^2 D)$
and the cumulative-release fraction at one day and at thirty
days using the long-time single-mode approximation.
\end{exercise}

\paragraph{Solution sketch.}
$L = 8 \times 10^{-4}$ cm. $L^2 = 6.4 \times 10^{-7}$ cm$^2$.
$\pi^2 D = 9.87 \times 2 \times 10^{-13} = 1.97 \times 10^{-12}$
cm$^2$/s. $\tau = 6.4 \times 10^{-7} / 1.97 \times 10^{-12}
\approx 3.2 \times 10^5$ s $\approx 3.8$ days. At one day,
$t/\tau = 1/3.8 = 0.26$, so $M_t/M_\infty \approx 1 - (8/\pi^2)
\exp(-0.26) \approx 1 - 0.81 \times 0.77 = 0.37$ (37 percent
released). At thirty days, $t/\tau = 7.9$, so $M_t/M_\infty
\approx 1 - 0.81 \times \exp(-7.9) \approx 1.0$ (effectively all
released). The cross-check using the full Crank series in
\texttt{code/drug\_elution\_diffusion.py} gives roughly 38 percent
at one day and ~99.97 percent at thirty days.

\begin{exercise}
A Bjork-Shiley convexo-concave valve at 70-degree opening angle in
the mitral position has an annual strut-fracture rate of
0.4 percent per patient-year. Compute (a) the probability that an
individual valve fractures within five years; (b) the expected
number of fractures over five years among a cohort of 1000
implanted patients; (c) the expected mortality from these
fractures, using the 70 percent single-leg-then-strut-fracture
case-fatality reported in \cite{paper:v6c12-blot-bjork-shiley-2005}.
Assume constant hazard.
\end{exercise}

\paragraph{Solution sketch.}
(a) With constant hazard $\lambda = 0.004$ per year, the survival
function is $S(t) = e^{-\lambda t}$, so $P(\text{fracture by 5 yr})
= 1 - e^{-5 \times 0.004} = 1 - e^{-0.02} = 1 - 0.9802 = 0.0198$,
or approximately 2 percent. (b) Expected fractures = $1000 \times
0.0198 \approx 20$ fractures over five years. (c) Expected
mortality = $20 \times 0.70 = 14$ deaths over five years among
the cohort. The calculation illustrates why the 29 mm and larger
mitral valves at 70-degree opening were withdrawn in 1983 and
why implanted-patient explantation was not the default
recommendation: the per-year hazard is small enough that the
cumulative five-year risk is modest, and elective explantation of
valve patients in their 60s and 70s has comparable surgical
mortality.

\begin{exercise}
The synthesised FDA recall data in
\texttt{data/fda\_recalls\_by\_class.csv} carry counts by class
(I/II/III) and by device category. Compute (a) the total recall
count; (b) the share of each class; (c) the device category with
the highest absolute count and the device category with the highest
share of Class I (most serious) recalls.
\end{exercise}

\paragraph{Solution sketch.}
Run \texttt{code/recall\_event\_classification.py}. (a) Total
recalls = $42 + 18 + \ldots + 8 = 1572$. (b) Class I = $42 + 18 +
9 + 15 + 21 + 7 + 6 + 11 = 129$ (8.2 percent). Class II = $1195$
(76.0 percent). Class III = $248$ (15.8 percent). (c) Surgical
has the highest absolute count (Class II = 221). Cardiovascular
has the highest Class I count (42), reflecting the proportion of
implantable cardiovascular devices in the recall stream.
Compare \cite{paper:v6c12-amato-recall-data-2014} for the
historical decadal trend.

\subsection*{Derivation}\pathtag{standard}

\begin{exercise}
Starting from the one-dimensional diffusion equation
$\partial C / \partial t = D \, \partial^2 C / \partial x^2$ with
uniform initial concentration $C_0$ in a slab $0 \le x \le L$,
zero-flux boundary at $x = 0$, and zero concentration at $x = L$,
derive the cumulative-release fraction $M_t / M_\infty$ as a
Fourier series. Identify the conditions under which the leading
term dominates.
\end{exercise}

\paragraph{Solution sketch.}
Separate variables: $C(x, t) = \sum_n A_n \cos((n + 1/2)\pi x / L)
\exp(-(n+1/2)^2 \pi^2 D t / L^2)$. Initial condition determines
$A_n$ by Fourier expansion of the constant $C_0$ on the cosine
basis: $A_n = (4 C_0 / \pi) (-1)^n / (2n + 1)$. Total released
mass $M_t = \int_0^L C_0 \, dx - \int_0^L C(x, t) \, dx$. The
infinite mass $M_\infty = C_0 L$. The integration gives
$M_t/M_\infty = 1 - (8/\pi^2) \sum_n (1/(2n+1)^2)
\exp(-(2n+1)^2 \pi^2 D t / L^2)$, which is Crank's solution
\cite{paper:v6c12-crank-diffusion-1975}. The leading term ($n = 0$)
dominates when $\pi^2 D t / L^2 \gtrsim 0.3$, that is, after
$t \gtrsim 0.03 \, L^2/D$. For $L = 10 \,\mu\text{m}$, $D = 10^{-13}
$ cm$^2$/s, this is approximately 30,000 seconds = 8 hours.

\begin{exercise}
Derive the equilibrium fibrous-capsule thickness predicted by the
simple model
\[
\frac{dh}{dt} = \alpha M(h) - \beta h,
\]
where $h$ is capsule thickness, $M(h)$ is the macrophage
recruitment rate (which decreases with $h$ as the capsule
isolates the implant), and $\beta h$ is collagen turnover.
Use $M(h) = M_0 / (1 + h/h_0)$ and find $h_\infty$ in closed form.
\end{exercise}

\paragraph{Solution sketch.}
At steady state $dh/dt = 0$, giving $\alpha M_0 / (1 + h_\infty/h_0)
= \beta h_\infty$. Solving: $\alpha M_0 = \beta h_\infty
(1 + h_\infty/h_0) = \beta h_\infty + (\beta/h_0) h_\infty^2$. This
is a quadratic in $h_\infty$: $(\beta/h_0) h_\infty^2 + \beta
h_\infty - \alpha M_0 = 0$. The positive root is
\[
h_\infty = \frac{-\beta + \sqrt{\beta^2 + 4 \alpha M_0 \beta / h_0}}
{2 \beta / h_0}
= \frac{h_0}{2}\left(-1 + \sqrt{1 + 4 \alpha M_0 / (\beta h_0)}\right).
\]
In the limit of weak recruitment ($\alpha M_0 \ll \beta h_0$),
$h_\infty \approx \alpha M_0 / \beta$. In the limit of strong
recruitment ($\alpha M_0 \gg \beta h_0$), $h_\infty \approx
\sqrt{\alpha M_0 h_0 / \beta}$. The square-root scaling is the
useful regime for thick capsules. The model is illustrative;
real capsule maturation involves multiple cell populations and is
better captured by the time-resolved simulation in
\texttt{code/host\_response\_timeline.py}.

\begin{exercise}
For a tilting-disc heart valve with two outlet-strut weld toes,
each modelled as an independent crack-initiation site with Paris-
Erdogan law $da/dN = C (\Delta K)^m$, $\Delta K = Y \Delta \sigma
\sqrt{\pi a}$, derive the expression for the cycle count to
critical crack length $a_c$. Show that the dependence on initial
flaw size $a_0$ is steep when $m > 2$.
\end{exercise}

\paragraph{Solution sketch.}
$da/dN = C (Y \Delta \sigma)^m (\pi a)^{m/2}$. Separating
variables: $da / a^{m/2} = C (Y \Delta \sigma)^m \pi^{m/2} dN$.
Integrating $a$ from $a_0$ to $a_c$ and $N$ from 0 to $N_f$:
\[
N_f = \frac{1}{C (Y \Delta \sigma)^m \pi^{m/2}}
\int_{a_0}^{a_c} a^{-m/2} \, da.
\]
For $m \ne 2$: $\int a^{-m/2} da = a^{1 - m/2} / (1 - m/2)$, so
\[
N_f = \frac{a_c^{1 - m/2} - a_0^{1 - m/2}}
{C (Y \Delta \sigma)^m \pi^{m/2} (1 - m/2)}.
\]
For $m > 2$, $1 - m/2 < 0$, so $a_0^{1 - m/2}$ is large and
dominates. Reducing $a_0$ by a factor of 2 increases $N_f$ by a
factor of $2^{m/2 - 1}$; for $m = 3.2$ (representative of
Haynes-25, the BSCC strut material), this is $2^{0.6} \approx 1.5$.
For $m = 4$, it is $2^1 = 2$. The implication: half-decade
improvements in initial-flaw quality control (welding,
electropolishing, surface finishing) can double or more the
fatigue life of the implanted component. This is the
quality-control failure the BSCC story turns on
\cite{paper:v6c12-paris-erdogan-1963,
paper:v6c12-vanderlugt-strut-fracture-1993,
web:v6c12-doj-shiley-1994-settlement}.

\subsection*{Estimation}\pathtag{core}

\begin{exercise}
Estimate the average number of cardiac cycles delivered to a
mechanical heart valve over a 25-year implant lifetime in a patient
with mean resting heart rate of 72 beats per minute and a 20
percent activity fraction at elevated heart rate (mean 110 bpm).
\end{exercise}

\paragraph{Solution sketch.}
Mean weighted rate = $0.80 \times 72 + 0.20 \times 110 = 57.6 +
22 = 79.6$ bpm. Cycles per year = $79.6 \times 60 \times 24 \times
365.25 \approx 4.2 \times 10^7$. Over 25 years, $1.05 \times 10^9$
cycles. The order of magnitude is $10^9$ cycles, consistent with
the 35 to 40 million cycles per year used in the BSCC registry
entry as a population figure. The estimate informs the fatigue-
design horizon for any cyclically loaded implant component.

\begin{exercise}
Estimate the residual ethylene oxide concentration in a polymeric
medical device after a standard EtO sterilisation cycle and a 14-day
aeration period, assuming exponential outgassing with characteristic
time of 4 days. Compare to the ISO 10993-7 24-hour patient-exposure
limit for a limited-contact device (4 mg / device for a single-use
product).
\end{exercise}

\paragraph{Solution sketch.}
Initial post-cycle residual is in the range $50$ to $500$ ppm
EtO in the bulk polymer for many materials
\cite{paper:v6c12-zhang-eo-residuals-2019}. Take 100 ppm as a
reference. For a 5-gram device: initial EtO mass $\approx 100 \times
10^{-6} \times 5 = 5 \times 10^{-4}$ g $= 0.5$ mg. After 14 days of
aeration with $\tau = 4$ days, the residual is $0.5 \times
\exp(-14/4) = 0.5 \times 0.030 = 0.015$ mg = 15 $\mu$g, far below
the 4 mg/device limit. The estimate shows that the aeration step
is essential (the un-aerated residual would already approach the
limit for a larger device) and that for typical devices the
standard 14-day aeration is sufficient with a comfortable margin.

\begin{exercise}
Estimate the metal-ion release rate from a metal-on-metal hip
implant operating under edge-loading conditions, assuming a wear
rate of 1 mm$^3$/year of CoCrMo alloy and a density of 8.3
g/cm$^3$. Express the daily cobalt release in micrograms, assuming
cobalt content of approximately 64 percent by mass.
\end{exercise}

\paragraph{Solution sketch.}
Annual mass loss = $1 \,\text{mm}^3 \times 8.3 \,\text{g/cm}^3
\times (1\,\text{cm}^3 / 10^3\,\text{mm}^3) = 8.3 \times 10^{-3}$
g/yr = 8.3 mg/yr. Cobalt fraction: $0.64 \times 8.3 = 5.3$ mg/yr.
Daily release: $5.3 \,\text{mg/yr} / 365 = 14.5 \,\mu\text{g/day}$.
For comparison, total dietary cobalt intake is approximately
12 $\mu$g/day (cobalt-deficient diet) to 100 $\mu$g/day (high
intake). The implant release is in the same order of magnitude as
dietary intake. Edge-loading conditions can increase the wear rate
by an order of magnitude to 10 to 30 mm$^3$/year; the corresponding
daily release approaches several hundred micrograms, which is
consistent with the elevated blood cobalt concentrations
documented in the ASR cohort \cite{paper:v6c12-cohen-mom-bmj-2012,
paper:v6c12-langton-asr-2011}.

\begin{exercise}
Estimate the lifetime risk that a pacemaker battery requires
replacement before the device's nominal end-of-service. Assume an
exponential battery lifetime with mean 8 years and that the
device's nominal service life is 10 years.
\end{exercise}

\paragraph{Solution sketch.}
$P(\text{battery fails before 10 yr}) = 1 - e^{-10/8} = 1 - e^{-1.25}
= 1 - 0.287 = 0.71$, or 71 percent. The number is high because the
exponential model assumes constant-hazard failure, which is wrong
for batteries (real pacemaker batteries have a near-zero failure rate
for most of their life followed by a sharp decline near end-of-
service). A more realistic Weibull model with shape parameter $\beta
= 3$ to $5$ and scale parameter set so that mean life is 8 years
gives a much lower probability of premature replacement. The
exercise's point: model selection matters for life-prediction. The
constant-hazard model is wrong for a battery, right for an
electronic-component random failure. Mixing the two requires care.

\subsection*{Simulation}\pathtag{mastery}

\begin{exercise}
Run \texttt{code/host\_response\_timeline.py} with default
parameters. Report the time at which the macrophage population
reaches its peak and the asymptotic fibrous-capsule density.
Modify the macrophage-fusion rate to be 2x the default and
re-run; explain the qualitative effect on the giant-cell trajectory.
\end{exercise}

\paragraph{Solution sketch.}
The default run shows neutrophils peaking around day 0.5 (acute
phase), declining sharply through day 2; macrophages rise through
days 5-15 and persist at a chronic level; giant cells form
slowly through days 10-30; the capsule density rises monotonically
toward the asymptote (90 percent of $C_{\max}$ reached around day
40 to 60). Doubling the fusion rate $k_{\text{fusion}}$ accelerates
giant-cell formation: the giant-cell trajectory rises earlier and
reaches higher asymptotic values. Because the giant-cell term
multiplies the macrophage-recruitment-driven collagen term, the
capsule also matures slightly faster. The exercise illustrates
why fibrosis-modulating biomaterial surface treatments aim to
suppress macrophage fusion rather than only macrophage recruitment.

\begin{exercise}
Modify \texttt{code/drug\_elution\_diffusion.py} to plot
$M_t/M_\infty$ for $L \in \{5, 10, 15, 20\}$ $\mu$m at $D = 10^{-13}$
cm$^2$/s. Identify the $L$ that gives $M_t/M_\infty = 0.5$ at $t =
30$ days. Compare to the leading-mode characteristic time $\tau
= L^2/(\pi^2 D)$.
\end{exercise}

\paragraph{Solution sketch.}
The half-release time is approximately $0.2 \, L^2 / D$ in the
short-time Crank regime. For $D = 10^{-13}$ cm$^2$/s and $t = 30
$ days $= 2.6 \times 10^6$ s, the required $L^2 = 5 t D = 1.3
\times 10^{-6}$ cm$^2$, so $L \approx 11.4 \,\mu\text{m}$. The
$\tau$ for $L = 11.4 \,\mu\text{m}$ is $(11.4 \times 10^{-4})^2 /
(\pi^2 \times 10^{-13}) = 1.3 \times 10^{6}/9.87 = 1.3 \times 10^5
$ s $\approx 1.5$ days, suggesting the leading-mode approximation
has crossed into the long-time regime by 30 days. The closed-form
solution from
\cite{paper:v6c12-crank-diffusion-1975} integrated in the code
gives the precise answer.

\begin{exercise}
Run \texttt{code/bjork\_shiley\_fracture\_paris.py}. Report the
predicted fracture rates per year for the 27 mm aortic at 60
degrees and the 29 mm mitral at 70 degrees. Compare to the
synthesised data in \texttt{data/bjork\_shiley\_fracture\_by\_size.csv}.
\end{exercise}

\paragraph{Solution sketch.}
The Paris-Erdogan integrator with $a_0 = 50 \,\mu\text{m}$ and
representative Haynes-25 parameters predicts the 27 mm aortic at
60 degrees at well below $10^{-4}$/year (consistent with the
$\approx 6 \times 10^{-4}$ in the data; the model under-predicts
because the initial flaw size 50 $\mu$m is a representative not
worst-case figure), and the 29 mm mitral at 70 degrees at a rate
two to three orders of magnitude higher (the model captures the
qualitative ordering). The exercise illustrates the Paris-
Erdogan law as a working tool while exposing the sensitivity to
$a_0$ and to the stress-amplitude heuristic. A useful follow-up
is to vary $a_0$ across the population: even a 50 to 100 $\mu$m
distribution of initial flaws produces the observed size-dependent
fracture-rate spread.

\begin{exercise}
Monte Carlo a five-step downstream manufacturing process for a
sterile single-use device whose step yields are $\{0.95, 0.92,
0.90, 0.95, 0.98\}$ with sterility-failure probabilities of
$\{10^{-7}, 10^{-7}, 10^{-7}, 10^{-7}, 10^{-7}\}$. Report (a) the
mean overall yield; (b) the mean expected SAL of the released
product; (c) the implications for parametric release.
\end{exercise}

\paragraph{Solution sketch.}
(a) Mean yield = product of step yields = $0.95 \times 0.92 \times
0.90 \times 0.95 \times 0.98 \approx 0.733$. (b) Combined sterility
failure $\approx 5 \times 10^{-7}$ assuming independence; the
released product's SAL is approximately $5 \times 10^{-7}$, better
than the standard $10^{-6}$ target. (c) Parametric release is
feasible because the combined sterility margin exceeds $10^{-6}$
by a comfortable factor; the more sensitive parameter is the
step-yield product, which moves with raw-material quality and
process drift. Monitoring step yields catches drift; monitoring
sterility margins catches calibration loss in the sterilisation
cycle.

\subsection*{Design}\pathtag{standard}

\begin{exercise}
Design a biological-evaluation programme for a novel implantable
glucose sensor (subcutaneous, intended for 90-day use). Specify
the ISO 10993 tests required, the test articles or extracts to be
used, the timepoints, and the justification for any tests
leveraged from a predicate device.
\end{exercise}

\paragraph{Discussion.}
A 90-day subcutaneous implant is implant/long-term under ISO
10993-1. Required tests: cytotoxicity (L929 extract, ISO 10993-5),
sensitisation (guinea-pig maximisation or LLNA), irritation,
acute systemic toxicity, subacute and chronic toxicity (with
weight-of-evidence on the 90-day endpoint), genotoxicity (Ames,
mouse micronucleus), implantation (ISO 10993-6 subcutaneous site
at 1, 4, 13 weeks), and carcinogenicity (extension or
weight-of-evidence to the 1-year limit). Test articles: the
finished, sterilised, packaged sensor; the polymer membrane;
sensor extract in saline and in cottonseed oil at 37 $^{\circ}$C
for 72 hours. Predicate-leverage: if a predicate sensor with the
same polymer membrane and same sterilisation has biological-
evaluation data on file, the cytotoxicity, sensitisation,
irritation, and acute systemic toxicity can be leveraged; the
implant- and chronic-endpoint data must be repeated for the new
geometry. See \cite{std:v6c12-iso-10993-1-2018,
std:v6c12-iso-10993-6-2016}.

\begin{exercise}
Design a sterilisation cycle for a device that combines stainless-
steel components, a UHMWPE bearing surface, and an embedded
lithium battery. State the cycle, the rationale for the choice,
and the validation programme.
\end{exercise}

\paragraph{Discussion.}
No single terminal sterilisation cycle is compatible with all
three components: autoclave damages UHMWPE and battery; gamma and
e-beam oxidise UHMWPE and degrade battery; EtO is compatible with
all three if the residuals can be managed. Choose EtO at 30 to 40
$^{\circ}$C, 600 mg/L, 4 hours, with 14-day aeration. Validate via
ISO 11135 bioburden-based dose-setting, including triplicate
fractional positive cycles to demonstrate microbial-population
inactivation kinetics, full cycles with biological indicators
($B.\,atrophaeus$ spore strips) placed at hardest-to-sterilise
locations, and EtO residual testing per ISO 10993-7 at one-day,
seven-day, and fourteen-day aeration timepoints. Compatibility of
the lithium battery with EtO is well-established and the cycle is
the canonical choice for combined electronic-polymer implants.
See \cite{std:v6c12-iso-11135-2014} and the matrix in
\texttt{code/sterilisation\_dose\_matrix.py}.

\begin{exercise}
Design the post-market surveillance programme for a new Class III
implantable cardioverter-defibrillator. Specify the data streams,
the trigger thresholds for regulatory action, the reporting
cadence, and the relationship to the device's clinical-trial
arm.
\end{exercise}

\paragraph{Discussion.}
Data streams: (i) device telemetry from the home-monitoring system
(battery state, lead impedance, shock events, sensing errors);
(ii) clinician-reported events (lead fracture, infection,
generator replacement); (iii) patient adverse-event reports
through MAUDE \cite{web:v6c12-fda-maude-2024}; (iv) post-approval
study cohort (PAS, mandated for Class III). Trigger thresholds:
shock-event rate exceeding 1 percent per device-year above the
clinical-trial 95 percent CI; lead-fracture rate exceeding 0.5
percent per device-year; infection rate exceeding 1 percent at one
year. Reporting cadence: quarterly to the FDA for the first two
years, annual thereafter; immediate (24-hour) reporting of any
death or serious injury. The PAS arm is the regulator's principal
post-approval evidence-generation mechanism and the device's
performance against the PAS cohort is the principal predictor of
labelling expansion or constriction over the device's commercial
life. See ISO 14971 \cite{std:v6c12-iso-14971-2019} for the
post-production information loop the surveillance programme
implements.

\subsection*{Diagnosis}\pathtag{core}

\begin{exercise}
A new total hip replacement system is cleared via 510(k) in 2024.
Two years later, the national joint registry reports an elevated
revision rate (4 percent at two years, against an expected 1 to 2
percent). Name three hypotheses for the elevated rate, the data
streams that would distinguish them, and the regulatory action
each would warrant.
\end{exercise}

\paragraph{Discussion.}
Hypothesis 1: device-specific design flaw (taper-junction
fretting in modular components, for example). Distinguish via
explanted-device analysis and per-surgeon revision-rate analysis.
Action: device recall and design correction if confirmed.
Hypothesis 2: surgeon learning curve (new instrumentation,
inadequate training). Distinguish via per-surgeon and per-centre
revision-rate distributions. Action: training programme,
restricted-use labelling. Hypothesis 3: patient-selection drift
(younger, more active patients than in clinical trial).
Distinguish via cohort demographics in registry. Action:
relabelling, prescriber education. The DePuy ASR case is the
canonical demonstration that the first hypothesis is rare in the
absolute and, when present, is the failure mode the registry is
best positioned to catch
\cite{paper:v6c12-smith-asr-revision-2012, web:v6c12-aoanjrr-2023}.

\begin{exercise}
A drug-eluting stent shows a late-thrombosis rate of 0.6 percent
per year, against the manufacturer's predicted 0.2 percent per
year. Name two engineering hypotheses for the discrepancy and the
in vitro / ex vivo / clinical tests that would distinguish them.
\end{exercise}

\paragraph{Discussion.}
Hypothesis 1: delayed endothelialisation due to the
antiproliferative drug suppressing endothelial recovery
\cite{paper:v6c12-finn-des-pathology-2007}. Distinguish via OCT
(optical coherence tomography) of strut coverage at 6 to 12
months in a clinical follow-up cohort. Hypothesis 2: residual
polymer-coating thrombogenicity from inflammatory response to
incompletely-elaborated coating. Distinguish via ex vivo
flow-loop blood-contact testing with heparinised blood and
platelet-deposition imaging. The two hypotheses are not exclusive
and historically both contributed to the first-generation DES
late-thrombosis problem; the second-generation DES platforms
addressed both via biocompatible polymers and thinner struts.

\begin{exercise}
A pacemaker manufacturer reports an unexplained cluster of
malfunction events at one clinical site. The events are not
reproducible at the manufacturer's bench. Name three hypotheses
the manufacturer should test before concluding the events are
operator error.
\end{exercise}

\paragraph{Discussion.}
Hypothesis 1: electromagnetic interference (EMI) specific to the
site (an MRI scanner nearby, an electrosurgical unit, a wireless
device population). Distinguish via on-site EMI survey, in-the-
field telemetry capture, and EMI lab reproduction of the site's
exposure profile. Hypothesis 2: lead-wire fracture in a specific
position or activity not captured in bench fatigue testing.
Distinguish via X-ray imaging of explanted leads from the affected
patients and fatigue analysis under the site's documented patient
activity profile. Hypothesis 3: software interaction with a
specific home-monitoring transceiver firmware version that the
clinical site uses. Distinguish via firmware regression testing
against the version-specific telemetry profile. The third
hypothesis is increasingly common as cardiovascular devices become
networked; the IEC 62304 software life-cycle process must address
deployed-environment variability.

\subsection*{Judgement}\pathtag{core}

\begin{exercise}
A small manufacturer wants to use the 510(k) pathway to introduce
a novel surgical-mesh material whose substantial equivalence rests
on a chain of three predicate devices, the earliest of which was
cleared in 1993. The 1993 predicate had its own predicate chain
that extends back to a pre-amendment device. Comment on the
manufacturer's strategy: what is the regulator likely to ask, and
what is the manufacturer's strongest defence?
\end{exercise}

\paragraph{Discussion.}
The regulator is likely to ask whether the predicate chain's
substantial-equivalence determinations carry adequate evidence for
the new material's biocompatibility and long-term implant
performance. The IOM 2011 report \cite{paper:v6c12-iom-510k-2011}
flagged exactly this pattern as a recurring weakness. The
manufacturer's strongest defence is direct biological-evaluation
data on the new material (full ISO 10993 programme), clinical-
literature evidence for the material in adjacent applications,
and a registry-based or claims-based comparison of the predicate
chain's actual long-term performance against expected baselines.
A weaker but commonly invoked defence is the predicate's regulatory
history of no recalls or adverse-event reports; this defence is
weaker because absence of recall is not evidence of safety,
particularly for low-volume long-implant-duration devices.
A De Novo pathway may be more defensible if the predicate chain
is genuinely strained.

\begin{exercise}
A manufacturer of a Class III implantable neurostimulator
identifies a software bug that, under specific stimulation
parameter combinations not approved in the labelling, could
deliver an over-stimulation event. The bug has not produced a
reported adverse event. State the manufacturer's regulatory
obligations and the considerations weighing for and against an
immediate field action.
\end{exercise}

\paragraph{Discussion.}
Under 21 CFR 803 the manufacturer must report the bug to the FDA
as a malfunction likely to cause or contribute to death or serious
injury if it were to recur, even if no event has been reported.
Under the QSR's CAPA requirements \cite{law:v6c12-21cfr-820} the
manufacturer must investigate the root cause, contain the bug,
and verify the correction. The question of field action depends on
the bug's reachability under labelled use (if unreachable, a
labelling update plus a clinician advisory may suffice; if
reachable, a software update via the home-monitoring system or a
physical service visit is required). The Therac-25 case is the
recurring example of why ``no reported event yet'' is not a
reliable signal: medical-device software events are under-reported,
particularly when the operator believes the event was an operator
error \cite{acc:therac-25-1985-1987, paper:leveson-turner1993,
std:v6c12-iec-62304-2015}. The manufacturer's risk-management file
(ISO 14971 \cite{std:v6c12-iso-14971-2019}) should drive the
decision via documented risk acceptance criteria, not via marketing
or commercial considerations.

\subsection*{Failure analysis}\pathtag{core}

\begin{exercise}
Trace the Bjork-Shiley convexo-concave failure chain from
manufacturing weld defect to patient mortality. Identify the
quality-control, regulatory, and clinical interventions that, at
each link, would have changed the outcome. Distinguish
interventions that were feasible in 1979 from those that became
feasible only after 1990.
\end{exercise}

\paragraph{Discussion.}
Chain: (i) weld defect at manufacturing; (ii) defective valve
released to production stream; (iii) PMA supplement pathway without
new clinical-outcome data; (iv) post-market reports of strut
fracture; (v) AECL-equivalent slow recognition of the recurring
fracture pattern; (vi) deferred recall on per-event-mortality
analysis; (vii) patient mortality on SLF event. Feasible-1979
interventions: tighter weld inspection at (i), production lot
rejection at (ii), additional clinical evidence requirement at
(iii). Feasible-after-1990 interventions: mandatory adverse-event
reporting under the 1990 SMDA at (iv); centralised joint or device
registry at (v); structured patient-notification programme at
(vi); see \cite{law:v6c12-safe-medical-devices-act-1990,
acc:bjork-shiley-cc-1979-1986}. The post-1990 interventions
together constitute the modern post-market device-surveillance
regime; the BSCC case is the canonical demonstration that they are
load-bearing safety functions, not bureaucratic overhead.

\begin{exercise}
A 32-year-old patient receives a metal-on-metal hip implant in
2008 (pre-recall). In 2011 (post-recall), the patient is
asymptomatic. Whole-blood cobalt is 9 $\mu$g/L. The MHRA threshold
for clinical investigation is 7 $\mu$g/L
\cite{web:v6c12-mhra-mom-2012}. Discuss the engineering, clinical,
and regulatory dimensions of the case. What is the appropriate
follow-up? Whose decision is it?
\end{exercise}

\paragraph{Discussion.}
Engineering dimension: the elevated cobalt is consistent with
wear-particle generation at the bearing, which is the predicted
failure-chain link from \cite{paper:v6c12-langton-asr-2011,
paper:v6c12-cohen-mom-bmj-2012}. The asymptomatic state does not
rule out ongoing wear and may not rule out incipient ALVAL or
pseudotumour formation. Clinical dimension: imaging (MARS MRI for
metal-artefact reduction sequences) is the standard follow-up for
elevated metal ions; serial measurements distinguish stable from
rising levels. Regulatory dimension: the MHRA threshold is for
investigation, not for explantation; the explantation decision
weighs the asymptomatic patient's surgical risk (revision arthroplasty
has approximately 1 percent perioperative mortality) against the
probability of progression to symptomatic pseudotumour or
loosening (perhaps 10 to 30 percent over five years for elevated-
cobalt patients on the available evidence). Whose decision: the
patient's, in consultation with their orthopaedic surgeon, with
the registry data and individual-patient risk-stratification as
inputs. The case illustrates that regulatory thresholds set
investigation triggers, not patient-management decisions; the
patient-management decision remains a clinical judgement informed
by the engineering of the device, the host's specific response,
and the patient's preferences.
